\providecommand{\Prim}{\operatorname{Prim}}
\providecommand{\Conn}{\operatorname{Conn}}

\section*{Local dictionary}
\addcontentsline{toc}{section}{Local dictionary}

\paragraph{\texorpdfstring{Coefficient category: Tate $\C$-vector spaces.}{Coefficient category: Tate C-vector spaces.}}
The coefficient-side linear algebra takes place in the symmetric
monoidal category $\mathrm{Tate}_\C$ of Tate $\C$-vector spaces:
ind-pro finite-dimensional $\C$-vector spaces equipped with linearly
compact lattices.  A linearly compact module is a lattice object, not
the whole Tate object.  The coefficient completed tensor product
$\widehat\otimes_{\mathrm{Tate}}$ is the topological tensor product on
$\mathrm{Tate}_\C$; continuous duality is the strict functor
$V\mapsto V^\vee_{\mathrm{cont}}=\Hom^{\mathrm{cts}}_\C(V,\C)$.  When
Dolbeault, de Rham, heat-kernel, or local-functional factors are
present, their analytic tensor product is the nuclear projective tensor
product \(\widehat\otimes_\pi\); the symbol \(\widehat\otimes\) below
means the completed product obtained from
\(\widehat\otimes_\pi\) on analytic factors and
\(\widehat\otimes_{\mathrm{Tate}}\) on coefficient factors.  No
bornological tensor product is used unless explicitly named.  The
$\C$-algebra $A=\C[[z_1,z_2]]$ is a complete linearly topologized
$\C$-algebra with the $\mathfrak m$-adic filtration
$\mathfrak m^k=(z_1,z_2)^kA$; $\mathfrak h=A/\C\!\cdot\!1$ is
pro-finite-dimensional in this filtration, and its strict continuous
dual is the Grothendieck--Matlis principal-part module
\[
  \mathfrak h^\vee_{\mathrm{cont}}
  =\mathcal P
  =\operatorname{Ann}_{\Res}(\C\cdot1)\subset H^2_{\mathfrak m}(A)\otimes\omega_A,
  \qquad
  \mathcal P=\bigoplus_{a+b>0}\C\,\rho_{a,b},
  \quad
  \rho_{a,b}=z_1^{-a-1}z_2^{-b-1}\,dz_1\,dz_2,
\]
ind-finite-dimensional, with the residue pairing
$\langle f,\omega\rangle=\operatorname{Res}_0\,f\omega$ a perfect pairing of Tate
$\C$-modules.  Inverse limits over Matlis windows
$\mathcal P_{\le N}=\bigoplus_{a+b\le N}\C\rho_{a,b}$ are read as
derived limits $\mathbf{R}\!\varprojlim_N\mathcal P_{\le N}$, with
$\varprojlim^1$-data tracked.  Polynomial calculations use the dense
ind-finite-dimensional sub-Lie algebra $\bar A=\C[z_1,z_2]/\C\cdot 1$;
completed CE/PV calculations use $\mathfrak h$ itself.  Factorization
algebras throughout are taken in the Costello--Gwilliam sense
(precosheaf on opens, Weiss-cover descent), refined to mixed-HT opens
on $\R^2\times\C^2$.  This is the analytic incarnation of the
factorisable-$\mathcal D$-module-on-$\mathrm{Ran}(\C^2)$ description ---
the two-complex-dimensional extension of the
$\mathcal D$-module-on-Ran presentation of curve chiral algebras
formalised in \cite{beilinson-drinfeld-chiral}*{\S 3}, equivalent to the analytic
form via the de Rham complex; the BV / heat-kernel formalism selects
Costello--Gwilliam as native.

\begin{defn}[Admissible host category]\label{defn:admissible-host-category}
  The category \(\mathcal B^{\mathrm{adm}}\) used in the formal-stalk
  theorem is the dg category whose objects are tuples
  \[
    (V,F^\bullet V,d,\widehat\otimes,V^\vee_{\mathrm{cts}},
    D_V,\mathrm{ev}_V,\mathrm{coev}_V,
    \underline{\mathrm{Hom}}^\bullet_{\mathrm{adm}},
    \operatorname{Roos},\,L_{\mathrm{sc}},\mathcal O_{\mathrm{adm}})
  \]
  where \((V,F^\bullet V,d)\) is a complete Hausdorff filtered cochain
  complex whose coefficient factor lies in \(\mathrm{Tate}_\C\) and
  whose analytic form or local-functional factors are nuclear
  locally convex complexes, \(V^\vee_{\mathrm{cts}}\) is the strict
  continuous dual, \(D_V\) is the declared coefficient dual used by BV
  pairings and kernels, \(L_{\mathrm{sc}}\) is the declared scalar line
  when one is present, and \(\mathcal O_{\mathrm{adm}}\) is the
  list of products, brackets, kernels, convolutions, trace transitions,
  and counterterm operations that are required to be continuous on
  \(V\).  Thus an object of
  \(\mathcal B^{\mathrm{adm}}\) is a complex together with its
  completed tensor product, internal Hom, strict and declared duality,
  Roos/strict-limit
  convention, scalar convention, and operation estimates.  It is not
  an unspecified background category.  The data satisfy the following
  conditions.
  \begin{enumerate}[label=\textup{(A\arabic*)}]
    \item \emph{Strict topology.}  The filtration is exhaustive and
      separated; differentials are strict continuous maps; cones and
      shifts carry the completed induced filtration.
    \item \emph{Completed tensor product.}  The monoidal product is the
      completed product \(\widehat\otimes\) obtained by applying the
      nuclear projective tensor product to analytic factors and the
      completed Tate tensor product to coefficient factors, then
      completing the algebraic tensor product for the convolution
      filtration
      \[
        F_n(V\widehat\otimes W)=
        \widehat{\sum_{i+j\leq n}\operatorname{im}
        (F_iV\otimes F_jW\to V\otimes W)}
      \]
      on the declared objects.  It preserves strict filtered
      quasi-isomorphisms in each variable.
    \item \emph{Continuous and declared duality.}  The strict continuous
      dual is
      \(V^\vee_{\mathrm{cts}}=\operatorname{Hom}^{\mathrm{cts}}_\C(V,\C)\)
      with the opposite Tate filtration.  The BV pairing in a theorem is
      carried by the declared coefficient dual \(D_V\), with continuous
      maps
      \[
        \mathrm{ev}_V\colon D_V\widehat\otimes V\to\C,\qquad
        \mathrm{coev}_V\colon \C\to V\widehat\otimes D_V .
      \]
      If \(D_V\ne V^\vee_{\mathrm{cts}}\), the theorem must state this
      coefficient enlargement and the topology in which
      \(\mathrm{coev}_V(1)\) converges.  No statement about \(D_V\) is
      then a statement about the strict continuous dual.
    \item \emph{Internal Hom.}  The dg Hom complex is
      \[
        \underline{\mathrm{Hom}}^k_{\mathcal B^{\mathrm{adm}}}(V,W)
        =
        \{\,f\colon V\to W[k]\mid f
        \text{ is continuous, filtered, and preserves the declared
        admissibility data}\,\},
      \]
      with differential
      \(d_{\mathrm{Hom}}f=d_W f-(-1)^k f d_V\).  The internal Hom
      used in convolution and centrality complexes is the closed
      subcomplex consisting of maps that preserve scalar lines, pairings,
      stable transition maps, and all operations listed in
      \(\mathcal O_{\mathrm{adm}}\).
    \item \emph{Limits.}  Inverse systems are evaluated by the Roos
      complex \(\operatorname{Roos}(V_\bullet)\), hence by
      \(\mathbf R\!\varprojlim V_n\).  A strict inverse limit may be used
      only when the Mittag--Leffler term
      \(\varprojlim\nolimits^1 H^{*-1}(V_n)\) is zero; otherwise that
      \(\varprojlim^1\)-class is an obstruction datum, not discarded.
    \item \emph{Scalar reduction.}  A Capelli or empty-trace scalar line
      \(L\simeq\C\) is a split central subobject.  The scalar-reduction
      functor is the quotient
      \[
        U_{\mathrm{sc}}(V,L)=V/L,
      \]
      while the unit of an algebra object is retained unless the theorem
      explicitly passes to Hamiltonian primitives.  This is the entire
      meaning of \(U_{\mathrm{sc}}\): it does not forget pro-Matlis
      completions, replace Roos limits by strict limits, discard
      Hochschild centrality homotopies, pass to stable traces, or erase a
      QME counterterm tower.
    \item \emph{Stable trace passage.}  For a tower of finite-rank
      invariant Koszul algebras \(\{\mathcal R_N^{GL_N}\}\), the stable
      trace functor is degreewise:
      \[
        \operatorname{StTr}_d(\mathcal R_\bullet)
        =
        \varprojlim_{N\geq d}
        \operatorname{Prim}H^0(\mathcal R_N^{GL_N},Q)_d
        \big/\C\cdot\operatorname{Tr}(1).
      \]
      Physical large-\(N\) states and cumulants are separate data.
      Here \(H^0(\mathcal R_N^{GL_N},Q)\) means the
      Chevalley--Eilenberg-length-zero invariant Koszul cohomology of
      the reduced polynomial brane complex.  Reductive exactness is
      Lemma~\ref{lem:reductive-invariants-koszul-exactness}; positive
      CE-length stabilizer classes of the quotient stack lie outside
      \(\operatorname{StTr}_d\).
    \item \emph{Operation estimates.}  Every product, bracket,
      coadjoint action, kernel contraction, Hochschild convolution, and
      Costello counterterm tower used on the object is accompanied by a
      continuity estimate in the chosen filtration.  A theorem invoking a
      diagonal Casimir or diagonal local-cohomology kernel must state the
      topology in which the displayed infinite sum converges.
  \end{enumerate}
  Morphisms in \(\mathcal B^{\mathrm{adm}}\) are degree-zero closed
  elements of the dg Hom complex.  Weak equivalences are strict
  filtered quasi-isomorphisms, or filtered quasi-isomorphisms after
  replacing the inverse system by its Roos complex.
\end{defn}

\begin{rmk}[Model objects and non-objects in \(\mathcal B^{\mathrm{adm}}\)]
\label{rmk:badm-models}
The finite-window object \(B_{\mathrm{fin}}\) is admissible: all sums
are finite, the tensor product is algebraic before completion, the
residue pairing is finite-dimensional, and the Schouten bracket,
Koszul differential, and trace maps are strict continuous maps.

For \(q>1\), let \(w(d)=q^d\).  The weighted Tate envelope
\(B_w\) of Definition~\ref{defn:wt-Bw-envelope} is admissible by
Theorem~\ref{thm:wt-canonical-envelope}.  Its declared coefficient
dual is
\(\mathfrak h^{\vee,w}_{\mathrm{cont}}\) of
Definition~\ref{defn:wt-coefficient-pair}; this is the product
coefficient module that carries the diagonal BV coevaluation.  It is
not the strict continuous dual of the Tate product
\(\mathfrak h_w\), whose strict continuous dual is the direct sum.
The product, Schouten bracket, Schouten differential, coadjoint
action, weighted diagonal Casimir, heat-kernel propagator, fixed-graph
renormalization-group contribution, and CE/PV map are the operations
listed in Theorem~\ref{thm:wt-canonical-envelope}.

The strict Tate product/direct-sum pair
\[
  \prod_{d\ge1}H_d,\qquad \bigoplus_{d\ge1}H_d^\vee
\]
has a continuous evaluation pairing, but it is not an admissible
coefficient object for the BV propagator used in the theorem: the
diagonal Casimir has nonzero component in every \(H_d^\vee\) and does
not lie in the completed tensor product with the direct-sum dual.
Consequently a theorem using the diagonal kernel in this strict pair
has no object of \(\mathcal B^{\mathrm{adm}}\) as its source.
\end{rmk}

\begin{prop}[Verified admissible finite Hamiltonian model]
\label{prop:verified-badm-finite-model}
  Let
  \[
    H_{\leq 2}
    =
    \operatorname{Span}_{\C}\{z_1,z_2,z_1^2,z_1z_2,z_2^2\}
    \subset \C[z_1,z_2]/\C\cdot1 .
  \]
  Since \(\{H_{\leq2},H_{\leq2}\}\subset H_{\leq2}\) modulo constants,
  \(H_{\leq2}\) is a finite-dimensional Hamiltonian Lie algebra.  Put
  \(P_{\leq2}=H_{\leq2}^{\vee}\) with the coadjoint action dual to this
  finite Lie bracket, and
  \(\mathfrak g_{\leq2}=H_{\leq2}\ltimes P_{\leq2}[1]\).  The completed
  finite CE/PV object
  \[
    B_{\leq2}
    =
    \widehat{\mathbf S}(H_{\leq2})
    \widehat\otimes
    \Lambda(H_{\leq2}^{\vee})
  \]
  with \(I\)-adic filtration \(F^rB_{\leq2}=I^r\), where
  \(I=(O_f,\theta^\lambda)\), differential
  \(d_\pi=[\frac12C^K_{IJ}O_K\theta^I\theta^J,-]_S\), ordinary completed
  tensor product over finite-dimensional coefficient spaces, continuous
  dual \(B_{\leq2}^{\vee,\mathrm{cont}}\), Roos convention equal to the
  constant inverse system, no scalar line \(L_{\mathrm{sc}}=0\), and
  operation list
  \[
    \mathcal O_{\leq2}
    =
    \{\text{multiplication},\,d_\pi,\,[\, ,\,]_S,\,
      H_{\leq2}\text{-coadjoint action},\,\text{finite Casimir}\}
  \]
  is an object of \(\mathcal B^{\mathrm{adm}}\).
\end{prop}

\begin{proof}
  The Poisson bracket of two generators of degree at most two has degree
  at most two, and the only degree-zero bracket is a constant, killed in
  \(\C[z_1,z_2]/\C\cdot1\).  Hence \(H_{\leq2}\) is a finite Lie algebra,
  and \(P_{\leq2}=H_{\leq2}^{\vee}\) is its ordinary finite coadjoint
  module.

  Conditions (A1)--(A8) of
  Definition~\ref{defn:admissible-host-category} are then checked inside
  finite-dimensional linear algebra.  For (A1), the \(I\)-adic completion
  is
  \[
    B_{\leq2}=\varprojlim_r B_{\leq2}/I^r,
  \]
  with finite-dimensional quotients; it is complete, Hausdorff, and
  separated.  Since \(d_\pi(O_K)=C^L_{KJ}O_L\theta^J\) and
  \(d_\pi(\theta^K)=-\frac12C^K_{IJ}\theta^I\theta^J\),
  \(d_\pi(I^r)\subset I^{r+1}\), so the differential is strict
  continuous.  For (A2), the tensor product is the \(I\)-adically
  completed tensor product of finite-dimensional quotients; over \(\C\)
  it preserves strict filtered quasi-isomorphisms.  For (A3), the
  continuous dual is the direct sum of the duals of the homogeneous
  finite-dimensional quotients, and the pairing with \(B_{\leq2}\) is
  degreewise perfect.  The finite diagonal
  \[
    \Delta_{\leq2}=\sum_i e_i\otimes e^i
    \in H_{\leq2}\otimes P_{\leq2}
  \]
  is an honest tensor, so no convergence hypothesis is hidden.

  For (A4), the internal Hom is the complex of continuous filtered maps
  preserving \(\mathcal O_{\leq2}\), with differential
  \(d_{\mathrm{Hom}}f=d f-(-1)^{|f|}f d\); preservation of a finite list
  of continuous operations cuts out a closed subcomplex.  For (A5), the
  inverse system is constant, so the Roos complex is the object itself
  and \(\varprojlim^1=0\).  For (A6), no Capelli or empty-trace scalar
  line is declared; the scalar-reduction functor is the identity and the
  unit of \(\widehat{\mathbf S}(H_{\leq2})\) is retained.  For (A7), no
  stable \(N\)-trace passage is part of \(\mathcal O_{\leq2}\); finite
  invariant trace towers are therefore not invoked.

  For (A8), all operation estimates are explicit in the \(I\)-adic
  filtration:
  \[
    I^rI^s\subset I^{r+s},\qquad
    d_\pi(I^r)\subset I^{r+1},\qquad
    [I^r,I^s]_S\subset I^{r+s-1}.
  \]
  The coadjoint action and the finite Casimir use only the finite pair
  \(H_{\leq2}\otimes P_{\leq2}\).  Thus every operation in
  \(\mathcal O_{\leq2}\) is continuous in the declared topology, and
  \(B_{\leq2}\) is a verified object of
  \(\mathcal B^{\mathrm{adm}}\).
\end{proof}

\begin{prop}[Verified admissible weighted Hamiltonian model]
\label{prop:verified-badm-weighted-model}
  Let \(w(d)=q^d\) with \(q>1\), let
  \((\mathfrak h_w,\mathfrak h^{\vee,w}_{\mathrm{cont}})\) be the
  weighted Tate coefficient pair of
  Definition~\ref{defn:wt-coefficient-pair}, and let \(B_w\) be the
  weighted Tate envelope of Definition~\ref{defn:wt-Bw-envelope}.  Give
  \(B_w\) the differential \(d_\pi\), its weighted completed tensor
  product, strict continuous dual \(B_w^{\vee,\mathrm{cts}}\), declared
  coefficient dual \(D_{B_w}\) induced from
  \(\mathfrak h^{\vee,w}_{\mathrm{cont}}\), Roos convention from the
  finite degree-window tower, scalar line \(L_{\mathrm{sc}}=0\), and
  operation list
  \[
    \mathcal O_w=
    \{\cdot,\ d_\pi,\ [-,-]_{\mathrm{Sch}},\
      \text{coadjoint action},\
      C_{\mathfrak h,w},\
      P^w_{\epsilon,L},\
      W_\Gamma(P^w_{\epsilon,L},I),\
      \Phi^w\}.
  \]
  Here \(C_{\mathfrak h,w}\) is the weighted diagonal Casimir,
  \(P^w_{\epsilon,L}\) is the weighted heat-kernel propagator, \(\Gamma\)
  ranges over fixed Costello graphs with polynomial-degree-bounded
  vertices, and \(\Phi^w\) is the weighted CE/PV map.  This tuple is an
  object of \(\mathcal B^{\mathrm{adm}}\).
\end{prop}

\begin{proof}
  Lemma~\ref{lem:wt-weights-exist} verifies the degree-weight
  conditions for \(w(d)=q^d\).  The tower of finite degree windows is
  strict Mittag--Leffler by Definition~\ref{defn:wt-degree-weight},
  hence inverse limits are computed by the Roos complex and have no
  hidden \(\varprojlim^1\)-term in the coevaluation tower.  The completed
  tensor product and declared coefficient dual are those of
  Definition~\ref{defn:wt-coefficient-pair}; the strict continuous dual
  \(B_w^{\vee,\mathrm{cts}}\) remains separate from \(D_{B_w}\), as
  required by Definition~\ref{defn:admissible-host-category}.

  Theorem~\ref{thm:wt-canonical-envelope} proves closure of \(B_w\)
  under multiplication, \(d_\pi\), and the Schouten bracket; convergence
  of \(C_{\mathfrak h,w}\) in
  \(B_w\widehat\otimes B_w\); extension of
  \(P^w_{\epsilon,L}\) as a continuous heat-kernel BV propagator;
  fixed-graph renormalization-group stability; and continuity of the
  cochain map \(\Phi^w\).  These are precisely the operation estimates
  required in (A8).  No scalar quotient or stable trace passage is part
  of this model, so (A6) and (A7) are identities.  Thus the data satisfy
  (A1)--(A8) of Definition~\ref{defn:admissible-host-category}.
\end{proof}

\begin{defn}[Dirac brane stack and boundary types]\label{defn:dirac-brane-stack-boundaries}
  Fix \(p\in\C^2_{\mathrm{hol}}\) and let
  \[
    L=\R_t\times\{0_s\}\times\{p\}\subset
    X=\R^2_{\mathrm{top}}\times\C^2_{\mathrm{hol}}
  \]
  be the worldline defect.  A Dirac brane at \(p\) is the formal
  skyscraper boundary condition supported on \(L\).  A stack of \(N\)
  Dirac branes means this boundary condition with Chan--Paton space
  \(\C^N\).  In a formal Darboux chart
  \(\widehat{\mathcal O}_{\C^2,p}\cong\C[[z_1,z_2]]\), the transverse
  coordinates act by
  \(z_i\mapsto\phi_i\in\mathfrak{gl}_N\), and the classical open
  observable chart is the quotient-stack Koszul cdga
  \[
    A^{\mathrm{cl}}_{\partial,N}
    =
    C^\bullet_{\mathrm{CE}}\!\left(\mathfrak{gl}_N,\,
      \mathrm{Sym}(\mathfrak{gl}_N\oplus\mathfrak{gl}_N
      \oplus\mathfrak{gl}_N[1])\right),
    \qquad
    Q\psi=[\phi_1,\phi_2].
  \]
  Equivalently, its transverse derived moduli chart is the derived
  commuting stack
  \[
    [\mu^{-1}_{\mathrm{der}}(0)/\mathrm{GL}_N],
    \qquad
    \mu(\phi_1,\phi_2)=[\phi_1,\phi_2].
  \]
  This definition is local and formal.  It is not the compact brane
  moduli problem and not the cyclic framed Hilbert-scheme branch.

  The Costello--Gwilliam brane-link boundary used in the local
  open--closed comparison is
  \[
    \partial_{\mathrm{link}}X
    =
    \partial(X\setminus\nu(L)),
  \]
  the boundary of an excised tubular neighbourhood of \(L\); in the
  local model it is the sphere bundle \(L\times S^4_\epsilon\).  The
  Kato--Nakayama boundary \(\partial X^{\mathrm{KN}}\) is a different
  global object: it is attached to a chosen log compactification
  \((\overline X,D)\).  A comparison between
  \(\partial_{\mathrm{link}}X\) and \(\partial X^{\mathrm{KN}}\) is
  extra matched-conventions data and is not part of the local
  formal-Darboux theorem.  In local brane-link constructions below,
  \(\partial X\) denotes \(\partial_{\mathrm{link}}X\); the
  Kato--Nakayama boundary is always written as \(\partial
  X^{\mathrm{KN}}\).
\end{defn}

\begin{defn}[Formal-Darboux stalk functor]\label{defn:formal-darboux-stalk-functor}
  For \(X=\R^k_{\mathrm{top}}\times\C^d_{\mathrm{hol}}\) and
  \(p\in\C^d\), the formal-Darboux stalk functor is the symmetric
  monoidal restriction functor
  \[
    \operatorname{FDSt}_p\colon
    \mathrm{ChirAlg}^{\mathrm{HT}}_{k,d}(X)
    \longrightarrow
    \mathrm{ChirAlg}^{\mathrm{HT}}_{k,d}
    (\R^k_{\mathrm{top}}\times\widehat{\C^d}_p),
    \qquad
    \mathcal F\longmapsto
    \mathcal F|_{\R^k_{\mathrm{top}}\times\widehat{\C^d}_p}.
  \]
  The target uses
  \(\widehat{\C^d}_p=\operatorname{Spf}\widehat{\mathcal O}_{\C^d,p}\)
  together with a formal Darboux coordinate system.  A different formal
  Darboux chart acts through
  \(\mathrm{Symp}^{\mathrm{form}}(\C^d)\); hence the stalk is canonical
  as a functor to the quotient by formal symplectomorphisms, and a
  chosen coordinate chart gives its displayed coordinate algebra.

  For the Hamiltonian BF bulk in this manuscript,
  \[
    \operatorname{FDSt}_p(\mathcal A^{\mathrm{cl}}_{\mathrm{bulk}})
    =
    C^\bullet_{\mathrm{CE},\mathrm{cont}}\!
    \left(\mathfrak h\ltimes\mathfrak h^\vee_{\mathrm{cont}}[1]\right),
    \qquad
    \mathfrak h=\C[[z_1,z_2]]/\C\cdot1.
  \]
  At \(N\) Dirac branes the boundary chart further applies
  \(z_i\mapsto\phi_i\in\mathfrak{gl}_N\) and the Koszul differential
  \(Q\psi=[\phi_1,\phi_2]\).  The trace-sector theorem computes the
  image of this functor after scalar reduction, stable trace passage,
  and admissible completion.  The full construction and equivariance
  proof are Definition~\ref{app:defn:Darboux-stalk-functor} and
  Proposition~\ref{app:prop:Darboux-equivariance}.
\end{defn}

\begin{defn}[Typed local source and target dictionary]
\label{defn:typed-local-source-target-dictionary}
The local theorem uses the following source and target objects.
\begin{enumerate}[label=\textup{(\roman*)}]
  \item \(\mathcal C^{\mathrm{op}}_\partial\) is the primitive boundary
    open factorisation category on the brane-link boundary.  A vacuum
    \(b\in\mathcal C^{\mathrm{op}}_\partial\) gives the boundary
    \(E_1\)-algebra
    \[
      A_b=\End_{\mathcal C^{\mathrm{op}}_\partial}(b).
    \]
    An unadorned \(A_\partial\), when used in a theorem-specific
    section, denotes such a chosen boundary \(E_1\)-algebra; the local
    main theorem uses the typed symbols \(A_b\) and
    \(A^{\mathrm{cl}}_{\partial,N}\).
  \item The formal coordinate algebra at \(b\) is
    \[
      A_b^{\mathrm{coord}}
      =
      \widehat{\mathcal O}_{\C^2,b}
      \cong \C[[z_1,z_2]].
    \]
    It contains Hamiltonians on the Darboux normal disk.  It is not the
    boundary algebra whose Hochschild cochains are computed.
  \item The Hamiltonian source is
    \[
      \mathfrak h=A_b^{\mathrm{coord}}/\C\cdot1,\qquad
      \bar A=\C[z_1,z_2]/\C\cdot1,\qquad
      P=\mathfrak h^\vee_{\mathrm{cont}}=\mathcal P,\qquad
      \mathfrak g=\mathfrak h\ltimes P[1].
    \]
    Here \(\bar A\) is the dense polynomial Hamiltonian subalgebra used
    for finite trace and sign calculations; completed
    Chevalley--Eilenberg/polyvector calculations use \(\mathfrak h\).
  \item The \(\mathfrak h\)-action on \(P\) is the residue coadjoint
    action, not the Matlis \(R\)-module multiplication:
    \[
      f\cdot\rho=\pi_{\mathrm{pp}}\{f,\rho\},\qquad
      z_1^p z_2^q\cdot\rho_{a,b}
      =-(pb-qa+p-q)\rho_{a-p+1,b-q+1},
    \]
    with negative-index targets zero.  This is the action used in the
    semidirect product \(\mathfrak g=\mathfrak h\ltimes P[1]\).  It
    extends continuously to
    \[
      \mathfrak h_{\mathfrak m\text{-adic}}\times
      \mathcal P_{\oplus}\longrightarrow\mathcal P_{\oplus},
    \]
    where \(\mathcal P_{\oplus}\) has the locally finite direct-sum
    topology: a fixed principal part sees only a finite jet of a formal
    Hamiltonian.  In the weighted pro-Matlis model the stronger
    estimate
    \[
      p^*_n(f\cdot\rho)\leq C_n\,p_{n+r}(f)\,p^*_{n+r}(\rho)
    \]
    is the continuity datum used by \(\mathcal B^{\mathrm{adm}}\).
    These statements are
    Proposition~\ref{prop:app-matlis-coadjoint-action},
    Lemma~\ref{lem:app-matlis-coadjoint-continuity}, and
    Proposition~\ref{prop:completion-promatlis-continuity}.
  \item At a stack of \(N\) Dirac branes, the finite-rank boundary
    algebra is
    \[
      A^{\mathrm{cl}}_{\partial,N}
      =\mathcal R_N
      =
      C^\bullet_{\mathrm{CE}}\!\left(\mathfrak{gl}_N,\,
        \mathrm{Sym}(\mathfrak{gl}_N\oplus\mathfrak{gl}_N
        \oplus\mathfrak{gl}_N[1])\right),
      \qquad Q\psi=[\phi_1,\phi_2].
    \]
    This is the quotient-stack Koszul cdga of the derived commuting
    stack.  Its augmentation ideal is
    \(\bar A^{\mathrm{cl}}_{\partial,N}\).  Fixed finite \(N\) retains
    trace relations, stabilizer cohomology, and unreduced centrality
    data.
  \item The stable scalar-reduced Hamiltonian trace algebra is
    \[
      A_{\partial,H}^{\mathrm{st}}
      :=
      \operatorname{TrSec}^{\mathrm{st}}_{\mathrm{adm}}(\mathcal R_\bullet)
      \cong
      \widehat{\mathbf S}_{\mathrm{Tate}}(\mathfrak h),
    \]
    where the construction first applies the scalar quotient
    \(U_{\mathrm{sc}}\), then degreewise stable trace passage
    \(\operatorname{StTr}\), then the Roos/strict-limit convention and
    the admissible completion of
    Definition~\ref{defn:admissible-host-category}.  This is a
    coordinate trace-sector algebra, not the whole fixed-\(N\)
    Hochschild centre.
  \item The finite-\(N\) Hochschild target is
    \[
      C^\bullet_{\mathrm{Hoch},\mathrm{adm}}
      (A^{\mathrm{cl}}_{\partial,N},A^{\mathrm{cl}}_{\partial,N}).
    \]
    The assignment \(c_f\mapsto\theta_f,\;u_f\mapsto J(f)\) is a
    generator boundary value into this complex until an admissible
    Hochschild-centre Maurer--Cartan lift and its centrality primitives
    have been supplied.
  \item The stable coordinate Hochschild target is
    \[
      HH^\bullet_{\mathrm{adm},\mathrm{HKR}}
      (A_{\partial,H}^{\mathrm{st}},A_{\partial,H}^{\mathrm{st}}),
    \]
    the Hochschild--Kostant--Rosenberg polyvector target of the
    stable scalar-reduced trace algebra.  It is the target of the
    formal CE/PV coordinate isomorphism.
  \item When \(C^\bullet_{\mathrm{ch}}(A_b,A_b)\) appears in a
    Deligne or chiral-centre statement, \(A_b\) denotes the boundary
    algebra with its restricted chiral/factorization structure.  The
    notation does not mean \(A_b^{\mathrm{coord}}\).  A global
    identification with the chiral derived \(E_2\)-centre requires the
    all-arity centre lift, descent, quantum master equation, anomaly,
    locality, and cone-contraction data recorded in
    Definition~\ref{defn:global-mixed-ht-deligne-obstruction-datum}.
\end{enumerate}
\end{defn}

\begin{defn}[Trace-sector reduction]\label{defn:trace-sector-reduction}
  Let
  \[
    \mathcal R_N=
    C^\bullet_{\mathrm{CE}}\bigl(\mathfrak{gl}_N,\,
       \mathrm{Sym}(\mathfrak{gl}_N\oplus\mathfrak{gl}_N
                    \oplus\mathfrak{gl}_N[1])\bigr),
    \qquad Q\psi=[\phi_1,\phi_2],
  \]
  and let \(T_N\) be a target dg algebra or Hochschild cochain complex
  in \(\mathcal B^{\mathrm{adm}}\).  Fix a polynomial window
  \(W\), a Chevalley--Eilenberg cochain-length bound \(L\), and
  \(\mathfrak h_{\leq W}=\mathfrak m/\mathfrak m^{W+1}\).  Put
  \[
    S_{W,L}:=
    C^{\leq L}_{\mathrm{CE}}
      \bigl(\mathfrak h_{\leq W}\ltimes
            \mathfrak h_{\leq W}^{\vee}[1]\bigr).
  \]
  A finite-window trace-sector datum on \(T_N\) is a continuous
  cochain map
  \[
    \Phi_{N,W,L}\colon
    S_{W,L}
    \longrightarrow T_N
  \]
  given by
  \[
    c_f\longmapsto\theta_f,\qquad
    u_f\longmapsto J(f),\qquad
    J(f)=\overline{\operatorname{Tr}f(\phi_1,\phi_2)},
  \]
  such that the scalar-reduced image is a strict admissible subcomplex
  of \(U_{\mathrm{sc}}T_N\).  Equivalently, the induced coimage map
  \[
    S_{W,L}/\ker(U_{\mathrm{sc}}\Phi_{N,W,L})
    \longrightarrow U_{\mathrm{sc}}T_N
  \]
  is a monomorphism in the chain category underlying
  \(\mathcal B^{\mathrm{adm}}\).  If this image is not strict, not
  closed under the differential, or not an admissible object of
  \(\mathcal B^{\mathrm{adm}}\), the window contributes an
  image-admissibility obstruction
  \(\operatorname{Ob}^{\mathrm{im}}_{N,W,L}\); no trace sector in
  \(\mathcal B^{\mathrm{adm}}\) is declared for that window.

  When this admissibility condition holds, the finite-window trace
  sector of \(T_N\) is the image/coimage subcomplex
  \[
    \operatorname{TrSec}_{W,L}(T_N)
    :=
    \operatorname{im}(U_{\mathrm{sc}}\Phi_{N,W,L})
    \cong
    S_{W,L}/\ker(U_{\mathrm{sc}}\Phi_{N,W,L})
    \subset U_{\mathrm{sc}}T_N .
  \]
  It is therefore a defined complex, not a list of trace generators.
  A bracketed or Hochschild trace sector further requires closure of
  this subcomplex under the restricted \(P_0\), cup, or Gerstenhaber
  operations used in the theorem; failure of this closure is recorded
  by the Hochschild-centrality obstruction complexes below.

  The stable admissible trace sector is
  \[
    \operatorname{TrSec}^{\mathrm{st}}_{\mathrm{adm}}(T_\bullet)
    :=
    \mathbf R\!\varprojlim_W\,
    \operatorname*{colim}_L\,
    \varprojlim_{N\geq \max(W,L+2)}
    \operatorname{TrSec}_{W,L}(T_N),
  \]
  with the Roos complex used for every non-strict inverse limit.
  Procesi--Razmyslov controls the polynomial trace generators in the
  range \(N\geq W\), and the \(L+2\) term is the Loday--Quillen--Tsygan
  finite cochain-window bound.

  For the coordinate polyvector target, the maps \(\Phi_{N,W,L}\) are
  the finite-window restrictions of the CE/PV coordinate isomorphism.
  For the Hochschild target
  \(T_N=C^\bullet_{\mathrm{Hoch},\mathrm{adm}}(\mathcal R_N,\mathcal R_N)\),
  such maps exist only after an admissible Hochschild-centre
  Maurer--Cartan lift has been chosen.  Without that lift,
  \(c_f\mapsto\theta_f,\;u_f\mapsto J(f)\) denotes a generator-level
  boundary value, not a Hochschild cochain map.

  We use three notations for the three resulting maps.
  \begin{enumerate}[label=\textup{(\roman*)}]
    \item \(\Phi_N^{\mathrm{gen}}\) denotes the dashed generator boundary
      value
      \[
        C^\bullet_{\mathrm{CE},\mathrm{cont}}
          (\widehat{\mathfrak g}_{\mathrm{cent}})
        \dashrightarrow
        C^\bullet_{\mathrm{Hoch},\mathrm{adm}}(\mathcal R_N,\mathcal R_N).
      \]
    \item \(\Phi_N^{\mathrm{Hoch}}\) denotes the Hochschild cochain map
      obtained from \(\Phi_N^{\mathrm{gen}}\) after an admissible
      Hochschild-centre Maurer--Cartan lift has been chosen.
    \item \(\Phi_{\mathrm{coord}}^{\mathrm{st}}\) denotes the
      scalar-stable Chevalley--Eilenberg/polyvector coordinate
      isomorphism
      \[
        C^\bullet_{\mathrm{CE},\mathrm{cont}}(\mathfrak g)
        \xrightarrow{\ \sim\ }
        \bigl(\mathrm{PV}_{\mathrm{adm}}
          (A_{\partial,H}^{\mathrm{st}}),d_\pi\bigr).
      \]
  \end{enumerate}

  The stable trace-sector construction is functorial on the following
  subcategory.  A morphism \(a_\bullet\colon T_\bullet\to T'_\bullet\)
  is allowed only when it is a morphism in \(\mathcal B^{\mathrm{adm}}\),
  preserves scalar lines, and carries each finite-window image
  \(\operatorname{im}(\Phi_{N,W,L})\) into
  \(\operatorname{im}(\Phi'_{N,W,L})\), either strictly or by a specified
  filtered homotopy compatible with the Roos transition maps.  Such a
  morphism induces
  \[
    \operatorname{TrSec}^{\mathrm{st}}_{\mathrm{adm}}(T_\bullet)
    \longrightarrow
    \operatorname{TrSec}^{\mathrm{st}}_{\mathrm{adm}}(T'_\bullet).
  \]

  The classical trace-coordinate reduction is the partial functor
  \[
    \operatorname{Red}_{\mathrm{cl,tr}}(T_\bullet)
    :=
    \mathrm{HKR}_{\mathrm{adm}}\!\left(
      \operatorname{TrSec}^{\mathrm{st}}_{\mathrm{adm}}(T_\bullet)
    \right)\big|_{\hbar=0},
  \]
  defined only on objects for which the Hochschild--Kostant--Rosenberg
  comparison, the Roos limit, the stable trace transition maps, and the
  scalar quotient above have been supplied.  If the Roos
  \(\varprojlim^1\)-class is nonzero, it remains part of
  \(\operatorname{Red}_{\mathrm{cl,tr}}\) as an obstruction component;
  it is not replaced by an ordinary inverse limit.

  The whole local bulk means
  \(C^\bullet_{\mathrm{CE},\mathrm{cont}}(\widehat{\mathfrak g}_{\mathrm{cent}})\)
  before applying \(U_{\mathrm{sc}}\), stable trace passage
  \(\operatorname{StTr}\), and \(\operatorname{TrSec}\).  The whole
  finite-\(N\) Hochschild centre means
  \(HH^\bullet(\mathcal R_N,\mathcal R_N)\).  Trace identities,
  stack-stabilizer cohomology, and unreduced centrality homotopies are
  outside \(\operatorname{TrSec}^{\mathrm{st}}_{\mathrm{adm}}\).
\end{defn}

\begin{defn}[Source complexes for the global mixed-HT Deligne obstruction]\label{defn:global-deligne-source-complexes}
  Put
  \[
    Z_\partial=
    Z^{\mathrm{der},\mathrm{Mor}}_{E_2^{\mathrm{HT}}}
    (\mathcal C^{\mathrm{op}}_\partial).
  \]
  A pre-open--closed comparison datum \(\mathrm{preOCA}_N\) is a
  filtered morphism
  \[
    \mathcal A^{\mathrm{cl}}_{\mathrm{bulk}}
    \longrightarrow
    \underline{\mathrm{End}}_{\mathrm{SC}^{\mathrm{HT}}_{(1,2)}}
    (\mathcal C^{\mathrm{op}}_\partial)
  \]
  in \(\mathcal B^{\mathrm{adm}}\), together with the
  formal-Darboux trace-sector fibre of Theorem~\ref{thm:main-local}.
  A centrality morphism to \(Z_\partial\) and the corresponding operadic
  homotopies promote this datum to \(\mathrm{OCA}_N\).

  The chiral \(E_2\) source complex is the homotopy fibre
  \[
    \mathfrak{Def}^{\mathrm{HT}}_{\mathrm{Mor}\mbox{-}E_2}
    (\mathcal C^{\mathrm{op}}_\partial,b_p)
    :=
    \operatorname{fib}\!\left[
    \mathrm{Conv}\bigl(\mathrm B\mathsf P,
      \operatorname{End}_{Z_\partial}\bigr)
    \xrightarrow{\operatorname{res}_{b_p}}
    \mathrm{Conv}\bigl(\mathrm B\mathsf Q,
      \operatorname{End}_{\operatorname{ev}_{b_p}(Z_\partial)^{\mathrm{tr}}}
      \bigr)\right][-1],
  \]
  where \(\mathsf P=\mathsf{hE}_2^{\mathrm{ch},\mathrm{HT}}\),
  \(\mathsf Q\) is its formal Hamiltonian-current restriction at
  \(b_p\), and \(\mathrm B\mathsf P,\mathrm B\mathsf Q\) are bar
  cooperads of chosen cofibrant dg-operad models.  The differential is
  twisted by the Maurer--Cartan element of the formal tree-kernel
  current action.  The full construction is
  Definition~\ref{app:defn:morita-chiral-E2-extension-complex}.

  For a collared Weiss cover \(\mathfrak U=\{U_a\}\) of
  \(\partial X\), write
  \(\mathcal A=\mathcal A^{\mathrm{cl}}_{\mathrm{bulk}}\).  The descent
  mapping complex is
  \[
    \mathcal E^\bullet_{\mathrm{desc}}(U_{a_0\cdots a_p})
    =
    \underline{\mathrm{Hom}}_{\mathcal B^{\mathrm{adm}}}^{\bullet}
    \bigl(
      \mathcal A|_{U_{a_0\cdots a_p}},
      Z_\partial|_{U_{a_0\cdots a_p}}\bigr),
  \]
  with differential
  \(d_{\mathcal E}\varphi=d_{Z_\partial}\varphi
  -(-1)^{|\varphi|}\varphi d_{\mathcal A}\), and
  \(\check C^\bullet_{\mathrm{Weiss}}(\mathfrak U;
  \mathcal E^\bullet_{\mathrm{desc}})\) has total differential
  \(\delta_{\check C}+(-1)^p d_{\mathcal E}\).
  This is Definition~\ref{app:defn:descent-obstruction-complex}.

  A Costello quantum master-equation source complex consists of a
  gauge-fixing operator, heat kernel \(K_t\), propagator
  \(P_{\varepsilon,L}\), scale-\(L\) BV Laplacian \(\Delta_L\), local
  interaction \(I[L]\), scalar-symbol projection
  \(\widehat\sigma_{\mathrm{sc},w}\), and admissible local-functional
  complex \(\mathfrak Q^\bullet_{w,\partial,\hbar}(I)\).  The
  differential is
  \[
    d_{\mathrm{QME}}^L
    =
    Q+\{I[L],-\}_{L}+\hbar\Delta_L.
  \]
  Put
  \[
    \mathrm{CE}_{\bar A}^{\mathbf 1}
    =
    C^\bullet_{\mathrm{Lie}}(\bar A;\C\gamma_{\mathbf 1}).
  \]
  The residual non-scalar complex is
  \[
    \mathcal Q^\bullet_w(I)
    =
    \ker\!\left(
    \widehat\sigma_{\mathrm{sc},w}\colon
    \mathfrak Q^\bullet_{w,\partial,\hbar}(I)
    \to \mathrm{CE}_{\bar A}^{\mathbf 1}[1]\right).
  \]
  Its inverse-window compatibility is recorded by the Roos
  \(\varprojlim^1H^0\)-class.

  The operadic anomaly and locality classes live in the
  \(\mathrm{SC}^{\mathrm{HT}}_{(1,2)}\)-convolution \(L_\infty\)-complex
  \[
    \mathrm{Conv}_{\mathrm{SC}}
    (\mathcal A^{\mathrm{cl}}_{\mathrm{bulk}},Z_\partial),
  \]
  the convolution complex for morphisms of
  \(\mathrm{SC}^{\mathrm{HT}}_{(1,2)}\)-algebras, twisted by the chosen
  source and target structures.  Its linear term is the Hom differential
  \(d_{Z_\partial}\varphi-(-1)^{|\varphi|}\varphi d_{\mathcal A}\), and
  its higher brackets are induced by the cooperadic decomposition.
\end{defn}

\begin{defn}[Global mixed-HT Deligne obstruction datum]\label{defn:global-mixed-ht-deligne-obstruction-datum}
  Let \(\mathrm{preOCA}_N\) be a pre-open--closed comparison datum in
  \(\mathcal B^{\mathrm{adm}}\) on the Costello--Gwilliam brane-link
  boundary \(\partial X\), with a fixed collared Weiss cover
  \(\mathfrak U\), source
  \(\mathcal A^{\mathrm{cl}}_{\mathrm{bulk}}\) and target
  \[
    Z_\partial=
    Z^{\mathrm{der},\mathrm{Mor}}_{E_2^{\mathrm{HT}}}
    (\mathcal C^{\mathrm{op}}_\partial).
  \]
  The global mixed-HT Deligne obstruction datum is the pair
  \[
    \mathfrak{Ob}_{\mathrm{Del}}^{\mathrm{glob}}
    =
    \left(
      \operatorname{Ob}^{\mathrm{global},5},
      \operatorname{Cone}(\mathrm{OCA}_N)
    \right),
  \]
  where
  \[
    \operatorname{Ob}^{\mathrm{global},5}
    =
    \Bigl(
      \operatorname{Ob}^{\mathrm{ch}\,E_2}_{\mathrm{Deligne}},
      \operatorname{Ob}^{\mathrm{desc}}_{\partial},
      \operatorname{Ob}^{\mathrm{QME}}_{\mathrm{all\mbox{-}order}},
      \operatorname{Ob}^{\mathrm{anom}},
      \operatorname{Ob}^{\mathrm{loc}}
    \Bigr).
  \]
  These entries are the following native classes:
  \[
  \begin{array}{rcl}
  \operatorname{Ob}^{\mathrm{ch}\,E_2}_{\mathrm{Deligne}}
  &\in&
  H^2\!\left(
  \mathfrak{Def}^{\mathrm{HT}}_{\mathrm{Mor}\mbox{-}E_2}
  (\mathcal C^{\mathrm{op}}_\partial,b_p)\right),
  \\[0.35em]
  \operatorname{Ob}^{\mathrm{desc}}_{\partial}
  &=&
  [\delta_{\check C}]\in
  \check H^2_{\mathrm{Weiss}}\!\left(
  \mathfrak U;H^{-1}(\mathcal E^\bullet_{\mathrm{desc}})\right),
  \\[0.35em]
  \operatorname{Ob}^{\mathrm{QME}}_{\mathrm{all\mbox{-}order}}
  &\in&
  \prod_{w,r} H^1\!\left(
  \ker\widehat\sigma_{\mathrm{sc},w},
  d_{\mathrm{QME}}^L\right)
  \times\varprojlim\nolimits^1 H^0,
  \\[0.35em]
  \operatorname{Ob}^{\mathrm{anom}}
  &\in&
  H^2\!\left(\mathrm{Conv}_{\mathrm{SC}}\!\left(
  \mathcal A^{\mathrm{cl}}_{\mathrm{bulk}},Z_\partial\right)\right),
  \\[0.35em]
  \operatorname{Ob}^{\mathrm{loc}}
  &\in&
  H^1\!\left(\partial X,
  \underline{\mathrm{Hom}}_{\mathcal B^{\mathrm{adm}}}^{0}
  (\mathcal A^{\mathrm{cl}}_{\mathrm{bulk}},Z_\partial)\right).
  \end{array}
  \]
  Here \(d_{\mathrm{QME}}^L=Q+\{I[L],-\}_{L}+\hbar\Delta_L\), defined
  only after the gauge fixing, propagator, finite-scale BV Laplacian,
  local interaction, and admissible counterterm complex have been fixed.
  Vanishing of \(\mathfrak{Ob}_{\mathrm{Del}}^{\mathrm{glob}}\) means
  specified primitives
  \[
    h_{E_2},\quad h_{\mathrm{desc}},\quad
    \{C_{w,r}\},\quad h_{\mathrm{anom}},\quad h_{\mathrm{loc}},\quad
    \kappa_{\operatorname{Cone}}
  \]
  satisfying the Maurer--Cartan lifting equation, the
  \v{C}ech--Weiss descent equation, the scale-\(L\) quantum master
  equation, the operadic anomaly equation, the factorization-locality
  equation, and
  \(d_{\operatorname{Cone}}\kappa_{\operatorname{Cone}}+
  \kappa_{\operatorname{Cone}}d_{\operatorname{Cone}}=
  \operatorname{id}_{\operatorname{Cone}}\).  A nonzero component, or a
  non-contractible comparison cone, obstructs the global mixed-HT
  Deligne quasi-isomorphism.  The formal-Darboux trace-sector theorem
  supplies only the restriction of \(\mathrm{preOCA}_N\) to
  \(\operatorname{TrSec}^{\mathrm{st}}_{\mathrm{adm}}\) and the Capelli
  projective scalar; it does not construct these primitives.
\end{defn}

\begin{ex}[Finite-window admissible model]\label{exa:finite-window-admissible-model}
  Fix \(W<\infty\), and let
  \(V_N=\mathfrak{gl}_N\phi_1\oplus\mathfrak{gl}_N\phi_2
  \oplus\mathfrak{gl}_N[1]\psi\).  Let
  \(I_N\subset\mathrm{Sym}(V_N)\) be the augmentation ideal for the
  total polynomial degree in \(\phi_1,\phi_2,\psi\).  The finite open
  Koszul window is
  \[
    \mathcal R_{N,\leq W}
    =
    C^\bullet_{\mathrm{CE}}\!\left(
      \mathfrak{gl}_N,\,
      \mathrm{Sym}(V_N)/I_N^{W+1}\right),
    \qquad
    Q\psi=[\phi_1,\phi_2]\bmod I_N^{W+1}.
  \]
  The derivation \(Q\) preserves \(I_N^{W+1}\), since it sends the
  degree-one generator \(\psi\) to the degree-two commutator.  The full
  subcategory on complexes generated by
  \[
    \mathfrak h_{\leq W}=\mathfrak m/\mathfrak m^{W+1},\qquad
    \mathcal P_{\leq W}=\bigoplus_{1\leq a+b\leq W}\C\,\rho_{a,b},
    \qquad
    \mathcal R_{N,\leq W}
  \]
  is an objectwise finite-dimensional model of
  \(\mathcal B^{\mathrm{adm}}\).  Tensor products are algebraic,
  continuous duality is ordinary linear duality, CE cochain length is
  bounded by the finite exterior algebra
  \(\bigwedge^\bullet\mathfrak{gl}_N^\vee\), and the Casimir
  \[
    C_{\leq W}=
    \sum_{1\leq a+b\leq W}z_1^az_2^b\otimes\rho_{a,b}
  \]
  is a finite sum.  For \(f\in\mathfrak h_{\leq W}\),
  \(J(f)=\overline{\operatorname{Tr}f(\phi_1,\phi_2)}\) lies in
  \(H^0(\mathcal R_{N,\leq W}^{GL_N},Q)\).  The adjacent-transposition
  homotopy
  \(Q\operatorname{Tr}(\psi vu)
    =\operatorname{Tr}([\phi_1,\phi_2]vu)\)
  remains inside \(\mathcal R_{N,\leq W}\) for words of length
  \(\leq W\).  Thus the finite-window trace-sector map is an honest
  finite-dimensional representative of the stable trace construction;
  its Hochschild version still requires the Hochschild-centre lift of
  Definition~\ref{defn:admissible-hochschild-centre-lift}.  The
  Procesi--Razmyslov stable trace identification in this window is used
  only for \(N\geq W\), or for \(N\geq\max(W,L+2)\) when a
  Chevalley--Eilenberg length \(L\) LQT projection is present.
\end{ex}

\begin{ex}[Non-admissible unweighted diagonal]\label{exa:unweighted-diagonal-nonadmissible}
  The naive infinite pair
  \[
    \mathfrak h_{\mathrm{naive}}=\prod_{a+b>0}\C\,z_1^az_2^b,\qquad
    \mathcal P_{\mathrm{alg}}=\bigoplus_{a+b>0}\C\,\rho_{a,b}
  \]
  with the product topology on \(\mathfrak h_{\mathrm{naive}}\) and the
  algebraic direct-sum topology on \(\mathcal P_{\mathrm{alg}}\) is not
  an admissible object for the local theorem.  The partial sums
  \[
    C_M=\sum_{1\leq a+b\leq M}z_1^az_2^b\otimes\rho_{a,b}
  \]
  have no limit in
  \(\mathfrak h_{\mathrm{naive}}\widehat\otimes\mathcal P_{\mathrm{alg}}\):
  their \(\mathcal P_{\mathrm{alg}}\)-support escapes every finite
  direct-sum stage.  Thus the formal diagonal Casimir is not a kernel in
  this topology.  Any theorem using the Casimir must use a finite window,
  a Roos/pro-Matlis target, or a weighted Tate completion that supplies
  the missing convergence.
\end{ex}

\begingroup
\small
\setlength{\parindent}{0pt}
\setlength{\tabcolsep}{3pt}
\renewcommand{\arraystretch}{1.15}

\[
  X=\R^2_{t,s}\times\C^2_{z_1,z_2},\qquad
  L=\R_t\times\{s=z_1=z_2=0\},\qquad
  \omega=dz_1\wedge dz_2.
\]

\paragraph{Unimodularity datum.}
\begin{longtable}{@{}>{\raggedright\arraybackslash}p{0.25\textwidth}
  >{\raggedright\arraybackslash}p{0.67\textwidth}@{}}
\hline
\textbf{Datum} & \textbf{Definition}\\
\hline
\endfirsthead
\hline
\textbf{Datum} & \textbf{Definition}\\
\hline
\endhead
\(\Omega_{\mathrm{loc}}\) &
The fixed holomorphic volume and symplectic form on the formal polydisk
\[
  \widehat{\C^2}_0=\operatorname{Spf}\C[[z_1,z_2]],
  \qquad
  \Omega_{\mathrm{loc}}=dz_1\wedge dz_2 .
\]
It trivializes the local canonical line and fixes the residue
orientation.\\
\hline
Poisson tensor and bracket &
The inverse bivector is
\[
  P=\partial_{z_1}\otimes\partial_{z_2}
    -\partial_{z_2}\otimes\partial_{z_1},
  \qquad
  \{f,g\}=\partial_{z_1}f\,\partial_{z_2}g
    -\partial_{z_2}f\,\partial_{z_1}g .
\]
Polynomial calculations use \(\bar A=\C[z_1,z_2]/\C\cdot1\); completed
CE/PV calculations use
\(\mathfrak h=\C[[z_1,z_2]]/\C\cdot1\).  Constants are removed because
they generate the zero Hamiltonian vector field.\\
\hline
Hamiltonian vector field &
The convention is
\[
  \iota_{X_f}\Omega_{\mathrm{loc}}=-df,\qquad
  X_f=(\partial_{z_1}f)\partial_{z_2}
      -(\partial_{z_2}f)\partial_{z_1}.
\]
Thus \(X_f(g)=\{f,g\}\) and
\([X_f,X_g]=X_{\{f,g\}}\) after projecting constants.\\
\hline
Unimodularity / divergence-free condition &
For a holomorphic vector field
\(X=a_1\partial_{z_1}+a_2\partial_{z_2}\),
\[
  \mathcal L_X\Omega_{\mathrm{loc}}
    =(\operatorname{div}_{\Omega_{\mathrm{loc}}}X)\Omega_{\mathrm{loc}},
  \qquad
  \operatorname{div}_{\Omega_{\mathrm{loc}}}X
    =\partial_{z_1}a_1+\partial_{z_2}a_2 .
\]
Every Hamiltonian vector field is divergence-free:
\[
  \operatorname{div}_{\Omega_{\mathrm{loc}}}X_f
  =-\partial_{z_1}\partial_{z_2}f
    +\partial_{z_2}\partial_{z_1}f=0 .
\]
Conversely, on the formal polydisk the holomorphic Poincare lemma identifies
\(\Ker(\partial_{\Omega_{\mathrm{loc}}}\colon\mathrm{PV}^1\to
\mathrm{PV}^0)\) with Hamiltonians modulo constants.\\
\hline
Closed BV pairing &
For the Hamiltonian BF fields
\[
  \Omega^\bullet(\R^2)\widehat\otimes\Omega^{0,\bullet}(\C^2)
  \widehat\otimes
  \bigl(\mathfrak h[1]\oplus
  \mathfrak h^\vee_{\mathrm{cont}}[2]\bigr),
\]
the BV pairing is the compactly supported local pairing obtained by
wedging the form factors, integrating against
\(\Omega_{\mathrm{loc}}\), and applying the continuous evaluation
\(\mathfrak h^\vee_{\mathrm{cont}}\otimes\mathfrak h\to\C\).  It has
degree \(-1\).  A Costello coefficient coevaluation kernel additionally
requires a Tate/weighted category in which the Casimir converges.\\
\hline
Residue / cyclic trace density &
The same orientation fixes the principal-part residue pairing
\[
  \rho_{a,b}=z_1^{-a-1}z_2^{-b-1}\,dz_1dz_2,\qquad
  \operatorname{Res}_0(\rho_{a,b}z_1^mz_2^n)
    =\delta_{a,m}\delta_{b,n},
\]
with \(a+b>0\) in the reduced cotangent module.  Hamiltonian vector
fields preserve \(dz_1dz_2\), so residue integration by parts gives the
coadjoint action.  On the open brane, the top Ext class is the residue
orientation dual to \(dz_1\wedge dz_2\), and the cyclic trace is
\[
  \operatorname{Tr}_{\mathrm{open}}(\eta\otimes u\otimes M)
  =\int_\R\eta\,[\varepsilon_1\varepsilon_2]u\,
    \operatorname{Tr}_N(M).
\]\\
\hline
\end{longtable}

\paragraph{Kinematic locality.}

\begin{thm}[Kinematic locality]
\label{thm:dict-kinematic-locality}
On the local pair \((X,L)\) with unimodularity datum
\(\omega=dz_1\wedge dz_2\) and reduced Hamiltonian Lie algebra
\(\mathfrak h=\C[[z_1,z_2]]/\C\cdot1\):
\begin{enumerate}
\item[(i)] The equal-time bracket kernel \(\delta(t-t')\) on two
  brane-time points, smeared against test functions \(a(t),b(t')\),
  gives the support-local \(P_0\)-bracket
  \(\{B_f(a),B_g(b)\}=B_{\{f,g\}}(ab)\) with \(ab\) the pointwise
  product on \(\R_t\); disjoint supports bracket to zero.
\item[(ii)] Every topological translation
  \(v\in T\R^2_{\mathrm{top}}\) is \(Q\)-exact, \(L_v=[Q,\iota_v]\),
  so \(H^\bullet(\mathrm{Obs}^{\mathrm{cl}})\) is locally constant
  along \(\R_t\) and inclusions of disks act by quasi-isomorphism.
\item[(iii)] Every antiholomorphic translation
  \(\partial/\partial\bar z_i\) is \(Q\)-exact,
  \(L_{\partial/\partial\bar z_i}=[Q,\iota_{\partial/\partial\bar z_i}]\),
  so the surviving \(\bar\partial\)-cohomology at the brane point is
  the formal jet algebra \(\mathfrak h\), Grothendieck-dually paired
  with
  \(\mathcal P=\operatorname{Ann}_{\mathrm{res}}(\C\cdot1)\subset
  H^2_{\mathfrak m}(\C[[z_1,z_2]])\,dz_1dz_2\).
\end{enumerate}
Items (ii) and (iii) hold on \(Q\)-cohomology; (i) is a chain-level
identity after distributional pairing.

Trace observables
\(\overline{\operatorname{Tr}\,f(\phi_1,\phi_2)}\) smeared as
\(B_f(a)\) satisfy (i)--(iii): support locality from the equal-time
bracket of the matrix-model fields, topological locality from the
\(Q\)-exact \(t\)-direction, formal-holomorphic locality from the
\(\bar\partial\)-exact transverse direction.  Shifted-cotangent
currents \(\Theta_\rho(B)\) satisfy (i) and (ii), and their
formal-holomorphic content is the residue current
\(\rho\in\mathcal P\) at the brane point.

The dg Lie algebra of an acceptable factorization algebra is therefore
locally constant in \(\R^2_{\mathrm{top}}\) and holomorphic in
\(B\subset\C^2_{\mathrm{hol}}\),
\[
  \Omega_c^\bullet(U)\widehat\otimes
  \Omega_c^{0,\bullet}(B)\widehat\otimes
  \mathfrak g,\qquad
  \mathfrak g=\mathfrak h\ltimes
  \mathfrak h^\vee_{\mathrm{cont}}[1],
\]
with \(P_0\)-bracket on smeared currents fixed by (i).  This produces
the \(\C^2\) holomorphic factorization algebra
\(\mathcal F_{\mathrm{Ham}}^{\mathrm{hol}}\) and its mixed
topological-current enhancement
\(\mathcal F_{\mathrm{Ham}}^{\mathrm{mix}}\).
\end{thm}

\paragraph{\texorpdfstring{Holomorphic factorization on \(\C^2\) and its curve
restrictions.}{Holomorphic factorization on 2 and its curve restrictions}}
The holomorphic object forced by
Theorem~\ref{thm:dict-kinematic-locality} is a
two-complex-dimensional factorization algebra on
\(\C^2_{\mathrm{hol}}\).  The taxonomy below records each object on
\(\R^2_{\mathrm{top}}\times\C^2_{\mathrm{hol}}\) and the restriction
data required to descend to a one-complex-dimensional factorization
algebra on a holomorphic curve.
\begin{longtable}{@{}>{\raggedright\arraybackslash}p{0.25\textwidth}
  >{\raggedright\arraybackslash}p{0.67\textwidth}@{}}
\hline
\textbf{Object} & \textbf{Definition and range}\\
\hline
\endfirsthead
\hline
\textbf{Object} & \textbf{Definition and range}\\
\hline
\endhead
Formal Hamiltonian holomorphic factorization algebra on \(\C^2\) &
The \(\C^2\) Hamiltonian holomorphic factorization algebra is
\[
  \mathcal F_{\mathrm{Ham}}^{\mathrm{hol}},
\]
the formal-Darboux Hamiltonian BF object on
\((\widehat{\C^2}_0,dz_1\wedge dz_2)\).  For a holomorphic polydisk
\(B\) its dg Lie algebra is
\[
  \Omega^{0,\bullet}_c(B)\widehat\otimes
  \mathfrak g,\qquad
  \mathfrak g=\mathfrak h\ltimes
  \mathfrak h^\vee_{\mathrm{cont}}[1],
  \quad
  \mathfrak h=\C[[z_1,z_2]]/\C\cdot1 ,
\]
with differential \(\bar\partial\) plus the Hamiltonian BF differential.
The formal stalk of observables is the completed CE/PV object
\[
  C^\bullet_{\mathrm{CE},\mathrm{cont}}(\mathfrak g)
  \simeq
  \mathrm{PV}_{\mathrm{coord}}(A_{\partial,\mathrm{coord}})
\]
in the coordinate/admissible sense recorded below.  This object is the
two-complex-dimensional holomorphic factorization algebra; a
one-complex-dimensional vertex algebra appears only after specifying a
curve restriction.\\
\hline
Native two-dimensional holomorphic singular product &
For two points \(z,w\in\C^2\), set \(\zeta_i=z_i-w_i\).  The singular
coefficient module is
\[
  H^2_\Delta(\C[[z,w]])\,d\zeta_1d\zeta_2
  =
  \bigoplus_{a,b\ge0}\C[[w_1,w_2]]\rho^\Delta_{a,b},
\]
where
\[
  \rho^\Delta_{a,b}
  =
  \zeta_1^{-a-1}\zeta_2^{-b-1}d\zeta_1d\zeta_2 .
\]
The normalized two-dimensional Cauchy kernel is
\[
  K^{(2)}_\Delta=\rho^\Delta_{0,0}
  =
  \frac{d\zeta_1d\zeta_2}{\zeta_1\zeta_2}.
\]
For \(x,y\in\mathfrak g\) the native current singular product is
\[
  \mathsf{J}_x(z)\mathsf{J}_y(w)\sim
  K^{(2)}_\Delta(z,w)\mathsf{J}_{[x,y]}(w),
\]
which gives
\(\mathsf{B}_f\mathsf{B}_g\sim K^{(2)}_\Delta\mathsf{B}_{\{f,g\}}\),
\(\mathsf{B}_f\mathsf{\Theta}_\rho\sim
K^{(2)}_\Delta\mathsf{\Theta}_{f\cdot\rho}\), and
\(\mathsf{\Theta}_\rho\mathsf{\Theta}_\sigma\sim0\).
Arity-\(n\) extension and operadic associativity:
Proposition~\ref{prop:diagonal-cohomology-operad}.
The quantum brane trace adds
\(\hbar N\,\omega(f,g)K^{(2)}_\Delta\mathbf1\), with
\(\omega(f,g)=[\{f,g\}]_0\).  This is the formal-coordinate
two-complex-dimensional OPE.  A one-variable Laurent expansion is the
additional curve-restriction datum recorded below.\\
\hline
Mixed topological-current enhancement over \(\R^2\) &
For a topological open \(U\subset\R^2_{\mathrm{top}}\) and holomorphic
polydisk \(B\subset\C^2_{\mathrm{hol}}\), the mixed current dg Lie algebra is
\[
  \mathfrak g^{\mathrm{mix}}_{U,B}
  =
  \Omega^\bullet_c(U)\widehat\otimes
  \Omega^{0,\bullet}_c(B)\widehat\otimes
  (\mathfrak h\ltimes\mathfrak h^\vee_{\mathrm{cont}}[1]).
\]
Its observables are locally constant in \(U\) and holomorphic in \(B\).
The brane-line current algebras below are its support-local
one-real-dimensional restrictions.\\
\hline
Weyl/Moyal open-string brane algebra and cyclic realisations &
The degree-zero quantum brane target is the Weyl/Moyal deformation of the
formal symplectic polydisk.  Put
\[
  P=\partial_{z_1}\otimes\partial_{z_2}
   -\partial_{z_2}\otimes\partial_{z_1}.
\]
Then
\[
  \begin{aligned}
  f*g&=m\circ\exp\!\left(\frac{\hbar}{2}P\right)(f\otimes g),\\
  \{f,g\}_\hbar&=\hbar^{-1}(f*g-g*f),
  \end{aligned}
\]
with constants quotiented in the Hamiltonian Lie algebra
\(\bar A_\hbar\).  At finite \(N\) the open model is the matrix Weyl
algebra generated by \((\phi_1)^i{}_j,(\phi_2)^k{}_l\) with
\[
  [(\phi_1)^i{}_j,(\phi_2)^k{}_l]
  =\hbar\,\delta^i_l\delta^k_j ,
\]
followed by the quantum moment-map reduction and the connected trace
projection.  Cyclic and supercyclic bar-cobar realisations are trace or
supertrace word complexes for matrix or supermatrix branes.\\
\hline
Defect/principal-part current algebra &
For an interval \(I\subset\R_{\mathrm{brane}}\),
\[
  \mathfrak h_I=\Omega^\bullet_c(I)\widehat\otimes\mathfrak h,\qquad
  \mathcal P_I=\mathcal D'_c(I)\widehat\otimes\mathcal P,\qquad
  \mathfrak g_I=\mathfrak h_I\ltimes\mathcal P_I[1],
\]
where
\[
  \mathcal P=\operatorname{Ann}_{\mathrm{res}}(\C\cdot1)
  \subset H^2_{\mathfrak m}(\C[[z_1,z_2]])\,dz_1dz_2 .
\]
The reduced target is
\[
  A^{\mathrm{pp}}_\partial(I)=
  \widehat{\mathbf S}(\Omega_c^0(I)\widehat\otimes\bar A)
  \widehat\otimes
  \widehat{\mathbf S}(\mathcal D'_c(I)\widehat\otimes\mathcal P[1]),
\]
with brackets
\[
  \begin{aligned}
  \{B_f(a),B_g(b)\}&=B_{\{f,g\}}(ab),\\
  \{B_f(a),\Theta_\rho(B)\}&=\Theta_{f\cdot\rho}(aB),\\
  \{\Theta_\rho(B),\Theta_\sigma(C)\}&=0 .
  \end{aligned}
\]
This is a support-local real-current \(P_0\) algebra with Matlis
principal parts in two holomorphic variables.\\
\hline
One-dimensional vertex/OPE restriction boundary &
A restriction of the \(\C^2\) Hamiltonian holomorphic
factorization algebra to a one-complex-dimensional vertex/OPE algebra
starts from the native diagonal singular product above and requires the extra datum
\[
  \mathfrak R_L
  =
  (\iota_L,\mathcal B_\perp,\mathsf R_{L,\mathcal B_\perp},
   V_L,\mathbf 1_L,T_L,Y_L,\zeta_{L,\hbar}),
\]
where
\(\iota_L\colon L=\operatorname{Spf}\C[[w]]
\hookrightarrow\widehat{\C^2}_0\) is a chosen holomorphic line or defect
germ, \(\mathcal B_\perp\) is a transverse boundary condition or normal
vacuum, and \(\mathsf R_{L,\mathcal B_\perp}\) is a restriction functor
producing a
one-complex-dimensional holomorphic factorization algebra
\[
  \mathcal F_{L,\mathcal B_\perp}(D)
  =
  \mathsf R_{L,\mathcal B_\perp}(\mathcal F_{\C^2})(D)
\]
for disks \(D\subset L\).  The functor must be built from a specified
tubular-neighbourhood, shriek/star restriction, or defect-coupling
construction, and must prove independence of the auxiliary transverse
radius, functoriality for disjoint disks, and translation covariance along
the coordinate \(w\).  The tuple also includes the restricted state
space \(V_L\), vacuum, translation, state-field map, and Zhu/zero-mode
map.  In the \(\C\times\R\) interface, such a \(V_L\) datum factors through
\(\mathcal R_{\C\times\R}\) and is identified with
\(V_{\mathrm{red}}\) in \(\mathfrak I_{\mathrm{ch}}\).\\
\hline
OPE data for a restricted curve object &
After such a restriction one must construct
\[
  V_L=H^\bullet(\mathcal F_{L,\mathcal B_\perp}(\Delta_0)),
  \qquad
  Y_L(a,w)b=\sum_{n\in\Z}(a_{(n)}b)\,w^{-n-1}.
\]
The factorization product for two disks centered at \(w\) and \(0\) must
have a finite-pole expansion
\[
  \mu_{w,0}(a\boxtimes b)
  \sim
  \sum_{n\ge0}(a_{(n)}b)\,w^{-n-1}
  +\mu^{\mathrm{reg}}_{w,0}(a,b),
\]
with \(a_{(n)}b=0\) for \(n\gg0\), a vacuum
\(\mathbf 1_L\), translation
\([T_L,Y_L(a,w)]=\partial_wY_L(a,w)\), and locality
\[
  (w_1-w_2)^N
  [Y_L(a,w_1),Y_L(b,w_2)]=0
\]
for \(N\gg0\).  These Laurent modes and finite singular orders belong
to the restricted curve object.  The two-complex-dimensional
holomorphic factorization algebra uses the singular coefficient
\(H^2_\Delta(\C[[z,w]])\,d\zeta_1d\zeta_2\); after curve restriction the
one-variable coefficient module is \(\C((w))/\C[[w]]\,dw\).\\
\hline
Compatibility with the Weyl/Moyal brane algebra &
A one-dimensional vertex/OPE restriction must also specify how the
degree-zero brane algebra acts on the restricted object.  Equivalently it
requires a continuous Zhu/zero-mode map
\[
  \zeta_{L,\hbar}\colon
  (\bar A_\hbar,*)\longrightarrow \operatorname{Zhu}(V_L)
\]
or fields \(J_f(w)\) satisfying
\[
  J_f(w)J_g(0)
  \sim
  \frac{\kappa_L(f,g)}{w^2}
  +
  \frac{J_{\{f,g\}_\hbar}(0)}{w},
\]
with the central level \(\kappa_L\) specified, and on the cotangent
current sector
\[
  [J_{f,(0)},\Theta_\rho]
  =
  \Theta_{\operatorname{ad}^{\vee}_{f,\hbar}\rho},
\]
together with
\(\zeta_{L,\hbar}(f*g)=\zeta_{L,\hbar}(f)\zeta_{L,\hbar}(g)\).
The Weyl/Moyal calculation gives \(f*g\) and the reduced
interval-current brackets; Theorem~\ref{thm:formal-two-dimensional-holomorphic-ope}
gives the native two-dimensional singular expansion.  The data
\(\zeta_{L,\hbar}\), the fields \(J_f(w)\), the central level, and the
finite Laurent curve expansion are separate one-dimensional restriction
constructions.\\
\hline
Chiral algebra interface tuple &
An externally constructed chiral algebra on a curve enters as the reduced
vertex object \(V_{\mathrm{red}}\) in
\[
  \mathfrak I_{\mathrm{ch}}
  =
  (\mathcal R_{\C\times\R},
   \mathcal F_{\mathrm{red}},
   V_{\mathrm{red}},
   Q_{\mathrm{BRST}},
   \mathbf 1,T,Y,
   \operatorname{wt},
   \zeta_{\hbar},
   H_{\mathrm{anom}}).
\]
Here \(\mathcal R_{\C\times\R}\) is a controlled pushforward to
\(\R_t\times\C_{z_1}\) retaining the \(z_2\)-mode or the full two-index
principal-part coefficient system
\[
  \mathfrak h_{\mathrm{red}}=\C[[z_1]][[z_2]]/\C\cdot1,\qquad
  \mathcal P_{\mathrm{red}}\subset
  H^2_{(z_1,z_2)}(\C[[z_1,z_2]])\,dz_1dz_2 .
\]
\(\mathcal F_{\mathrm{red}}\) is the pushed-forward factorization
algebra on \(\R_t\times\C_{z_1}\).  The interface must prove the
pushed-forward bracket, BV pairing, brane image, factorization-to-vertex
comparison, BRST nilpotence, Zhu/zero-mode map, Capelli normalization,
Moyal multiplicativity, the completed stable Hochschild/HKR/CE-PV
comparison, and anomaly transport.  The obstruction vector is
\[
  \operatorname{Ob}_{\mathfrak I_{\mathrm{ch}}}
  =
  (\operatorname{Ob}_{\mathrm{red}},
   \operatorname{Ob}_{\mathrm{vert}},
   \operatorname{Ob}_{\mathrm{Zhu}},
   \operatorname{Ob}_{\mathrm{nat}}).
\]
\(V_{\mathrm{red}}\) is a reduced shadow of
\(\mathcal F_{\mathrm{Ham}}^{\mathrm{hol}}\) on a holomorphic curve; the
two-complex-dimensional algebra remains
\(\mathcal F_{\mathrm{Ham}}^{\mathrm{hol}}\).\\
\hline
\end{longtable}

\paragraph{\texorpdfstring{Normal \(\Omega\)-background.}{Normal -background}}
The brane-preserving equivariant deformation scales the normal bundle
\(N_LX=\R_s\oplus\C_{z_1}\oplus\C_{z_2}\) with \(t\) fixed; the fixed
locus condition singles out this normal scaling.  The
objects and identities below define the deformation in the mixed de
Rham/Dolbeault coefficient category: the Cartan differential
\(Q_\Omega=Q+\iota_{V_\Omega}\), the localized normal contraction, the
equivariant CE/PV bracket, and the residue-versus-Euler normalization.
\begin{longtable}{@{}>{\raggedright\arraybackslash}p{0.25\textwidth}
  >{\raggedright\arraybackslash}p{0.67\textwidth}@{}}
\hline
\textbf{Object or convention} & \textbf{Definition}\\
\hline
\endfirsthead
\hline
\textbf{Object or convention} & \textbf{Definition}\\
\hline
\endhead
Ambient pair and fixed stratum &
The local pair is
\[
  X=\R^2_{t,s}\times\C^2_{z_1,z_2},\qquad
  L=\R_t\times\{s=z_1=z_2=0\}.
\]
The normal bundle of the brane line is
\[
  N_LX=\R_s\oplus\C_{z_1}\oplus\C_{z_2}.
\]
The \(\Omega\)-background scales this bundle with \(t\) fixed.\\
\hline
Equivariant parameters &
Write
\[
  T_\Omega
  =\C^*_{\epsilon_s}\times
   \C^*_{\epsilon_1}\times
   \C^*_{\epsilon_2}.
\]
The parameters \(\epsilon_s,\epsilon_1,\epsilon_2\) are equivariant
weights.  The open trace orientation uses the separate Ext generators
\(\varepsilon_1,\varepsilon_2\).
The action is
\[
  (s,z_1,z_2)\mapsto
  (\lambda_s s,\lambda_1z_1,\lambda_2z_2),
  \qquad t\mapsto t .
\]
Its infinitesimal vector field is
\[
  V_\Omega
  =
  \epsilon_s s\partial_s
  +\epsilon_1 z_1\partial_{z_1}
  +\epsilon_2 z_2\partial_{z_2}.
\]\\
\hline
Equivariant differential &
The Cartan differential on the local mixed complex is
\[
  Q_\Omega=Q+\iota_{V_\Omega},\qquad
  Q_\Omega^2=L_{V_\Omega}.
\]
Thus \(Q_\Omega\) is an honest differential on the
\(T_\Omega\)-invariant subcomplex.  Away from invariants the correct
object is the \(T_\Omega\)-equivariant curved Cartan complex with
curvature \(L_{V_\Omega}\).\\
\hline
Normal-weight localization ring &
For a finite \((N,M)\)-window let \(\mathsf X_{N,M}\) be the finite set
of nonzero moving characters that occur in that window.  It contains the
normal Hamiltonian and residue-dual characters
\[
  n\epsilon_s+a\epsilon_1+b\epsilon_2,\qquad
  n\epsilon_s-a\epsilon_1-b\epsilon_2
\]
whenever the corresponding moving summands occur.  The localized
coefficient ring is
\[
  R_\Omega^{N,M}
  =
  \C[\epsilon_s,\epsilon_1,\epsilon_2,\hbar_\omega]
  [\chi^{-1}\mid \chi\in\mathsf X_{N,M}].
\]
Equivalently, one inverts only the nonzero weights present in the finite
normal window.  This is the algebra in which moving normal
modes may be contracted.  The factor
\(\epsilon_s\epsilon_1\epsilon_2\) is inverted in the Euler-localization
branch after that branch has been declared.  The self-dual branch is a
specialization with its own allowed inversions: base-change first, then
invert only the surviving nonzero characters, leaving
\(\epsilon_1+\epsilon_2\) as a zero character on the self-dual branch.\\
\hline
Normal contraction datum &
A localized deformation retract
\[
  \mathsf C_\Omega=(\pi_L,j_L,h_\Omega)
  \colon
  \mathcal E_\Omega(X)
  \rightleftarrows
  \mathcal E_\Omega(L)
\]
for the mixed de Rham/Dolbeault coefficient complex, with
\[
  \pi_Lj_L=\operatorname{id},
  \qquad
  Q_\Omega h_\Omega+h_\Omega Q_\Omega
  =\operatorname{id}-j_L\pi_L
\]
on the localized invariant normal complex.  On a homogeneous normal summand of nonzero weight \(\lambda\),
\(h_\Omega\) is the Euler/Koszul homotopy divided by \(\lambda\), with
signs fixed by the de Rham degree in the \(s\)-direction and the
Dolbeault degree in the \(z_i\)-directions.\\
\hline
Hamiltonian and cotangent weights &
For \(H_{a,b}=z_1^az_2^b\), \(a+b>0\),
\[
  w(H_{a,b})=a\epsilon_1+b\epsilon_2.
\]
For the principal-part residue basis,
\[
  \rho_{a,b}=z_1^{-a-1}z_2^{-b-1}\,dz_1dz_2,
  \qquad
  w(\rho_{a,b})=-a\epsilon_1-b\epsilon_2 .
\]
The holomorphic symplectic character is
\[
  \chi_\omega=\epsilon_1+\epsilon_2,
  \qquad
  w(dz_1\wedge dz_2)=\chi_\omega,
  \qquad
  w(P)=-\chi_\omega .
\]
The Weyl/Moyal deformation parameter, when present, has
\[
  w(\hbar_{\mathrm{W}})=\chi_\omega.
\]\\
\hline
Three \(\hbar\)-parameters &
The Omega convention distinguishes three formal parameters:
\[
  \hbar_{\mathrm{QME}},\qquad
  \hbar_{\mathrm{W}},\qquad
  \hbar_\omega .
\]
Here \(\hbar_{\mathrm{QME}}\) is the Costello/BV loop-counting parameter
in graph and curvature equations; \(\hbar_{\mathrm{W}}\) is the
Weyl/Moyal deformation parameter in the open brane algebra; and
	\(\hbar_\omega\) is the equivariant line parameter of weight
	\(\chi_\omega\) used to scalarize the Poisson bracket.  A Weyl branch may
	identify \(\hbar_{\mathrm{W}}=\hbar_\omega\).  Identification of
	\(\hbar_{\mathrm{QME}}\) with another parameter is an explicit theorem
	convention.\\
\hline
Self-dual bracket &
On the self-dual locus
\[
  \epsilon_1+\epsilon_2=0
\]
the Poisson tensor has weight \(0\).  The Hamiltonian bracket
\(\{f,g\}\) is then a scalar \(T_\Omega\)-equivariant Lie bracket on
\(\mathfrak h\) after quotienting constants.\\
\hline
Refined bracket off the self-dual locus &
Off the self-dual locus the ordinary Poisson bracket has weight
\[
  w(\{f,g\})=w(f)+w(g)-(\epsilon_1+\epsilon_2).
\]
It is therefore a bracket valued in the inverse symplectic-character
line.  The scalar refined bracket is obtained only after adjoining the
line parameter \(\hbar_\omega\) of weight
\(w(\hbar_\omega)=\epsilon_1+\epsilon_2\):
\[
  [f,g]_\Omega=\hbar_\omega\{f,g\}.
\]
In the Weyl/Moyal specialization
\(\hbar_\omega=\hbar_{\mathrm{W}}\).  Before this scalarization, the
refined Hamiltonian algebra is valued in the inverse
symplectic-character line; after adjoining \(\hbar_\omega\) it becomes a
scalar equivariant Hamiltonian Lie algebra.\\
\hline
Equivariant CE/PV generator rule &
The equivariant CE/PV map keeps the same generator labels
\[
  c_{a,b}\mapsto\theta_{a,b},\qquad
  u_{a,b}\mapsto O_{a,b},
\]
but all differentials and brackets are read in the \(T_\Omega\)-graded
completed coordinate algebra.  In the refined convention the contraction
bracket and the \(d_\pi\)-terms carry the factor \(\hbar_\omega\).  In the
self-dual convention this factor has weight \(0\) and may be suppressed.\\
\hline
Equivariant quantum brane algebra &
The degree-zero open quantum brane target is the projective
central-character reduction
\[
  A^q_{N,\Omega}
  =
  \left(
  \operatorname{Weyl}_{\hbar_{\mathrm{W}}}(\operatorname{Mat}_N^2)/
  \mathcal W_N\widehat\mu_{\chi_N}(\mathfrak{sl}_N)
  \right)^{SL_N},
  \qquad
  YX-XY+\hbar_{\mathrm{W}}NI\equiv0
  \pmod{\mathcal W_N\widehat\mu(\mathfrak{sl}_N)},
\]
with
\[
  w(X)=\epsilon_1,\qquad
  w(Y)=\epsilon_2,\qquad
  w(\hbar_{\mathrm{W}})=\epsilon_1+\epsilon_2.
\]
The central-character relation is weight homogeneous.  Its scalar \(U(1)\)
contact term is the same Capelli line as in the non-equivariant
dictionary, now graded by \(\epsilon_1+\epsilon_2\).\\
\hline
Residue normalization &
The residue convention uses the formal residue orientation
\[
  \operatorname{Res}_{z_1=z_2=0}
  \bigl(z_1^{-1}z_2^{-1}dz_1dz_2\bigr)=1
\]
and the open top Ext coefficient
\([\varepsilon_1\varepsilon_2]\).  Its normalization constant is
\[
  \nu_\Omega^{\mathrm{res}}=1.
\]
In this convention the brane trace is
normalized by the residue/cyclic density.  Any
\(\epsilon_s\epsilon_1\epsilon_2\) factor is written explicitly as a
comparison constant.\\
\hline
Euler-localization normalization &
The smooth equivariant localization convention pushes from \(X\) to \(L\)
with the inverse normal Euler class
\[
  e_{T_\Omega}(N_LX)^{-1}
  =
  (\epsilon_s\epsilon_1\epsilon_2)^{-1},
\]
up to the orientation sign \(\sigma_s\) fixed by the chosen real
\(s\)-density and the complex orientation \(dz_1\wedge dz_2\).  Its
normalization constant is
\[
  \nu_\Omega^{\mathrm{Eu}}
  =
  \sigma_s(\epsilon_s\epsilon_1\epsilon_2)^{-1}.
\]
A theorem comparing this
convention with the residue convention must state the sign and the exact
factor.  Until then the two normalizations are separate branches.\\
\hline
Stratified factorization datum &
The brane-localized object is a stratified mixed HT factorization datum
on the pair \((X,L)\),
\[
  \mathcal F^{\mathrm{str}}_\Omega(X,L)
  =
  \bigl(\mathcal F_{\Omega,X\setminus L},
        \mathcal F_{\Omega,L},
        \mathcal M_{\Omega,L},
        j^*,i^!,\mu_{\mathrm{str}}\bigr),
\]
where the closed bulk stratum is the \(T_\Omega\)-equivariant mixed
bulk factorization algebra, the open defect stratum is the brane-line
current object, and \(\mathcal M_{\Omega,L}\) is the defect module.  Protected
observables become stratified factorization-homology traces in the
presence of the brane-defect QME and the trace-state normalization.\\
\hline
Open obstruction vector &
The residual obligations are
\[
  \begin{aligned}
  \operatorname{Ob}_\Omega
  &=
  \bigl(\operatorname{ob}_{\Omega,\mathrm{Cartan}},
        \operatorname{ob}_{\Omega,\mathrm{contr}},\\
  &\qquad
        \operatorname{ob}_{\Omega,\mathrm{CE/PV}},
        \operatorname{ob}_{\Omega,\mathrm{norm}},\\
  &\qquad
        \operatorname{ob}_{\Omega,\mathrm{strat}},
        \operatorname{ob}_{\Omega,\mathrm{QME}},\\
  &\qquad
        \operatorname{ob}_{\Omega,N\to\infty}\bigr).
  \end{aligned}
\]
These name, respectively, the mixed Cartan model, the localized normal
contraction, the equivariant CE/PV bracket comparison, the
residue-versus-Euler normalization theorem, the stratified factorization
structure, the Costello brane-defect graph/QME complex, and the physical
large-\(N\) trace-state theorem.\\
\hline
\end{longtable}

\paragraph{Canonical maps and obstruction symbols.}

\begin{longtable}{@{}>{\raggedright\arraybackslash}p{0.25\textwidth}
  >{\raggedright\arraybackslash}p{0.67\textwidth}@{}}
\hline
\textbf{Symbol} & \textbf{Meaning}\\
\hline
\endfirsthead
\hline
\textbf{Symbol} & \textbf{Meaning}\\
\hline
\endhead
\(\Phi_{\mathfrak l}\) &
Finite-dimensional cotangent CE/PV map
\[
  C^\bullet_{\mathrm{CE}}
  (\mathfrak l\ltimes \mathfrak l^\vee[1])
  \to \mathrm{PV}(\mathbf S(\mathfrak l)),
  \qquad c\mapsto\theta,\quad u\mapsto O .
\]\\
\hline
\(\Phi_{\mathrm{coord}}\) &
Formal-polydisk coordinate CE/PV map reconstructed from admissible finite
coordinate tests and pro-window coordinate cochains for
\(\mathfrak h=\C[[z_1,z_2]]/\C\cdot1\) and
\(\mathfrak g=\mathfrak h\ltimes
\mathfrak h^\vee_{\mathrm{cont}}[1]\).  It is a completed coordinate dg
commutative algebra isomorphism before any \(P_0\) bracket is asserted.\\
\hline
\(\Phi^B_{P_0}\) &
Restriction of \(\Phi_{\mathrm{coord}}\) to a bracket-admissible
subalgebra \(B\subset
\mathrm{PV}_{\mathrm{coord}}(A_{\partial,\mathrm{coord}})\):
\[
  \Phi^B_{P_0}\colon\Phi_{\mathrm{coord}}^{-1}(B)\xrightarrow{\sim}B .
\]
The completed \(P_0\) comparison in \(\mathcal B^{\mathrm{adm}}\)
(Theorem~\ref{thm:main-local}, axioms~(i)--(iv)): \(B\) carries a
complete Hausdorff topology, continuous inclusion in the coordinate
product, continuous multiplication, continuous
\(d_\pi^{\mathrm{coord}}\), and a continuous closed contraction bracket;
the criterion is Köthe bracket continuity
(Remark~\ref{rmk:consistency-probes}, check~(6)).
Boundary Hamiltonians inside this \(P_0\) comparison require the relevant
\(O_f\) and \(d_\pi^{\mathrm{coord}}O_f\) to lie in \(B\).\\
\hline
\(\Theta_B\) &
Kernel-admissible Maurer--Cartan tensor in a specified completed
convolution dg Lie algebra \(\mathcal K_B\),
\[
  \Theta_B=\sum_I H_I\otimes\theta^I+\sum_I\eta^I\otimes O_I ,
\]
attached to \(B\).  Asserted only when the displayed diagonal sums
converge and the convolution differential and bracket are continuous
in \(\mathcal K_B\); the kernel-admissibility test is the Maurer--Cartan
identity \(d\Theta_B+\tfrac12[\Theta_B,\Theta_B]=0\)
(Theorem~\ref{thm:main-local}, axiom~(v);
Remark~\ref{rmk:consistency-probes}, check~(7)).\\
\hline
\(\eta^I\) &
Shifted cotangent generator
\(\eta^I=(H_I)^\vee[1]\in\mathfrak h^\vee_{\mathrm{cont}}[1]\), for
\(I=(a,b)\) and \(a+b>0\).  It is the source of the \(O_I\) coordinate in
the CE/PV map.\\
\hline
\(\mathfrak U_{\mathfrak h}\) &
Admissible formal-local interface datum in the category
\(\mathcal B^{\mathrm{adm}}\) of
Definition~\ref{defn:admissible-host-category}.  The categorical universal
object is
\[
  \mathfrak U_{\mathfrak h}
  =(\Phi_{\mathrm{coord}},\{\Phi^B_{P_0}\}_B,\{\Theta_B\}_B),
\]
where \(B\) ranges over bracket-admissible subalgebras in the middle term
and kernel-admissible data in the last term.\\
\hline
\(B_{\mathrm{fin}},E_q\) &
Model admissibility tests.  \(B_{\mathrm{fin}}\) is the finite-support
coordinate subalgebra on a finite window.  \(E_q\), for \(q>1\), is the
weighted Köthe test target with seminorms
\[
  p_n(x)=\sup_{a,b}|x_{a,b}|q^{a+b}(1+a+b)^n .
\]
It is bracket-admissible only after the finite-word tensor convention
and the corresponding convolution estimate are fixed; it is
kernel-admissible only if the diagonal tensor converges.\\
\hline
\(\mathfrak o^{\mathrm{br}},
  \mathfrak o^{\mathrm{coad}}\) &
Transition obstruction terms for bracket and coadjoint compatibility in
the pro/mixed coordinate CE/PV comparison.\\
\hline
\(\mathfrak o^{\mathrm{Cas}}\) &
Strict product/direct-sum kernel obstruction for the Casimir,
propagator, and Costello endpoint problem.  It belongs to the endpoint
coefficient problem; the coordinate CE/PV identity is the separate
coordinate comparison.\\
\hline
\(\mathfrak o_{\mathrm{BV,end}}\) &
Endpoint obstruction tuple
\((\mathfrak o_{\mathrm{Cas}},
  \mathfrak o_{\mathrm{def}},
  \mathfrak o_{\mathrm{RG}},
  \mathfrak o_{\mathrm{QME}})\) controlling a Costello-compatible
coefficient endpoint.\\
\hline
\(\operatorname{Ob}_{\mathfrak C_{\mathrm{tar}}}\) &
Matched-conventions descent vector
\((\operatorname{per}_{\mathrm{tar}},
  \operatorname{ob}_{\mathrm{WR}},
  \operatorname{ob}_{\mathrm{anom}},
  \operatorname{ob}_{\mathrm{QME}},
  \operatorname{ob}_{\mathrm{conv}},
  \operatorname{ob}_{\mathrm{hyp}},
  \operatorname{ob}_{6r},
  \operatorname{ob}_{\mathrm{MD}})\),
indexed by the comparison target \(\mathfrak C_{\mathrm{tar}}\).\\
\hline
\end{longtable}

\paragraph{Non-scalar Costello/QME labels.}
The local obstruction complex named in
Problem~\ref{prob:quantum-p0-operation-realization} is recorded
below.  The scalar
projection \(\sigma_{\mathrm{sc}}\) is associated-graded data; the
projection \(\widehat\sigma_{\mathrm{sc}}\) is a filtered chain map
existing exactly when the successive scalar-lift obstructions vanish;
the reduced principal-part brackets live after de Rham-current
contraction; the unreduced non-scalar QME is the curvature equation
\(\operatorname{Curv}(K_{\hbar,w,I})=0\) of the bulk-to-defect kernel.
\begin{longtable}{@{}>{\raggedright\arraybackslash}p{0.25\textwidth}
  >{\raggedright\arraybackslash}p{0.67\textwidth}@{}}
\hline
\textbf{Symbol or datum} & \textbf{Meaning}\\
\hline
\endfirsthead
\hline
\textbf{Symbol or datum} & \textbf{Meaning}\\
\hline
\endhead
\(\mathfrak Q^\bullet_{w,\partial,\hbar}(I)\) &
Filtered brane-defect local-functional complex
\[
  \varprojlim_M
  \left(
    \mathcal O_{\mathrm{loc}}^{\mathrm{bd}}\bigl(
      \mathcal E_w^M|_I;
      A_{\partial,\hbar}^{\mathrm{pp},w,M}(I)\bigr),
    d_M
  \right),
  \qquad d_M=Q+\{I^w_{0,M},-\}_{\hbar,M}.
\]
The inverse limit is taken over finite Hamiltonian windows with the
\(\hbar\)-adic, scalar-contact, and finite-window filtrations.  The
	differential has degree \(+1\).  This is the target complex for the
	bulk-to-defect kernel and QME obstruction; the reduced current algebra is
	\(A_{\partial,\hbar}^{\mathrm{pp},w}(I)\).\\
\hline
\(F^\bullet\mathfrak Q^\bullet_{w,\partial,\hbar}(I)\) &
Complete decreasing closed scalar-contact filtration preserved by
\(d_{\mathrm{QME}}\).  In the ordinary scalar-reduced branch take
\(F^r=J_{\mathrm{sc}}^r\), where \(J_{\mathrm{sc}}\) is the closed ideal
of functionals with zero associated-graded scalar-contact symbol.  In the
balanced branch it is the closed filtration generated by
identity-endomorphism contractions inside the balanced scalar-contact
complex.  Its role is to define the associated-graded scalar symbol and
the obstruction tower for lifting that symbol.\\
\hline
\(\sigma_{\mathrm{sc}}\) &
Associated-graded scalar-symbol extractor
\[
  \sigma_{\mathrm{sc}}\colon
  \operatorname{gr}_F\mathfrak Q^\bullet_{w,\partial,\hbar}(I)
  \longrightarrow
  C^\bullet_{\mathrm{Lie}}(\bar A;\C\gamma_{\mathbf 1})[1][[\hbar]] .
\]
It has cochain degree \(0\) as an associated-graded map.  A projection of
the filtered QME complex is the lifted map
\(\widehat\sigma_{\mathrm{sc}}\).  On a degree-one scalar contact with two
Hamiltonian arguments it records the Lie two-cocycle coefficient of
\(\gamma_{\mathbf 1}\).  In the ordinary branch this leading symbol is
\(\hbar N[\bar c]\).\\
\hline
\(\widehat\sigma_{\mathrm{sc}}\) &
Actual scalar-contact projection: a continuous filtered degree-\(0\)
cochain map
\[
  \widehat\sigma_{\mathrm{sc}}\colon
  \mathfrak Q^\bullet_{w,\partial,\hbar}(I)
  \to
  C^\bullet_{\mathrm{Lie}}(\bar A;\C\gamma_{\mathbf 1})[1][[\hbar]]
\]
whose associated graded is \(\sigma_{\mathrm{sc}}\).  It exists exactly
when the scalar-lift obstruction classes
\(o_{\sigma,w}^{(r)}(I)\) vanish successively with compatible
finite-window choices.  The residual non-scalar complex is defined only
after this map has been chosen.  On a balanced scalar-contact complex
the proved projection is the zero chain map
\(\widehat\sigma_{\mathrm{bal,sc}}=0\).\\
\hline
\(o_{\sigma,w}^{(r)}(I)\) &
Successive scalar-contact lift obstruction.  For a filtered linear lift
\(s\) of \(\sigma_{\mathrm{sc}}\), the chain defect
\(\delta_s=d_{\mathrm{CE}}s-sd_{\mathrm{QME}}\) lowers \(F^\bullet\).
Its first nonzero part gives
\[
  o_{\sigma,w}^{(r)}(I)\in
  H^1\!\left(
    \operatorname{Hom}\bigl(
      \operatorname{gr}_F\mathfrak Q^\bullet_{w,\partial,\hbar}(I),
      C^\bullet_{\mathrm{Lie}}(\bar A;\C\gamma_{\mathbf 1})[1][[\hbar]]
    \bigr)_{-r}
  \right).
\]
In the ordinary scalar-reduced \(\lie{gl}_N\) branch the first
two-argument evaluation is
\(\hbar N\,\omega(f,g)\gamma_{\mathbf 1}\), nonzero on
\((f,g)=(z_1,z_2)\).\\
\hline
\(\mathfrak g^{w,\mathrm{cur}}_{\hbar,I}\) &
Current source for the unreduced non-scalar lift,
\[
  \bigl(\Omega^\bullet_c(I)\widehat\otimes\mathfrak h_{w,\hbar}\bigr)
  \ltimes
  \bigl(\mathcal D'_c(I)\widehat\otimes
  \mathfrak h^{\vee,w}_{\hbar,\mathrm{cont}}\bigr)[1].
\]
The topology is the completed current topology with the weighted
finite-window and \(\hbar\)-adic completions.  On CE coordinates,
\(c_{B,\rho}\) has degree \(1\) and
\(u_{a(t)dt\otimes f}\) has degree \(0\).  Before scalar reduction one
may replace the Hamiltonian source by the central-character source.\\
\hline
\(\kappa_{\hbar,w,I}\) &
Continuous bulk-to-defect kernel, equivalently a degree-\(0\)
Maurer--Cartan element in the completed convolution Lie algebra,
\[
  \kappa_{\hbar,w,I}\colon
  C^\bullet_{\mathrm{CE},\mathrm{cont}}
  (\mathfrak g^{w,\mathrm{cur}}_{\hbar,I})
  \longrightarrow
  \mathfrak Q^\bullet_{w,\partial,\hbar}(I),
\]
with generator values
\[
  u_{a(t)dt\otimes f}\mapsto B_f^{\mathrm{BV}}(a),
  \qquad
  c_{B,\rho}\mapsto\Theta_\rho^{\mathrm{BV}}(B).
\]
Its reduced image must recover \(B_f(a)\) and \(\Theta_\rho(B)\), but the
kernel itself lives before scalar reduction, before de Rham-current
contraction, and before replacing unreduced open observables by freely
adjoined current generators.\\
\hline
\(\operatorname{Ob}_{w,\partial,\hbar}\) &
Curvature cocycle of the bulk-to-defect kernel,
\[
  \operatorname{Ob}_{w,\partial,\hbar}
    =
  d_{\mathrm{QME}}\kappa_{\hbar,w,I}
  +\frac12[\kappa_{\hbar,w,I},\kappa_{\hbar,w,I}]_\hbar .
\]
It has QME degree \(1\).  The scalar branch first requires
\(\widehat\sigma_{\mathrm{sc}}(\operatorname{Ob}_{w,\partial,\hbar})=0\).
Only then does its residual class lie in the non-scalar complex.\\
\hline
\(\operatorname{Ob}^{\mathrm{red}}_{w,\partial,\hbar}\) &
Residual non-scalar obstruction after the scalar-contact projection has
been constructed and the scalar component has been killed:
\[
  [\operatorname{Ob}^{\mathrm{red}}_{w,\partial,\hbar}]
  \in
  H^1\!\left(
    \ker\widehat\sigma_{\mathrm{sc}},
    d_{\mathrm{QME}}
  \right).
\]
This class is defined after choosing the scalar-contact projection; the
reduced principal-part bracket formulas leave it as a residual
non-scalar obstruction.\\
\hline
\(\mathcal Q^\bullet_{w,M}(I)\) &
Finite-window residual complex
\[
  \mathcal Q^\bullet_{w,M}(I)
    =
  \ker(\widehat\sigma_{\mathrm{sc},M})
  \subset
  \left(
    \mathcal O_{\mathrm{loc}}^{\mathrm{bd}}\bigl(
      \mathcal E_w^M|_I;
      A_{\partial,\hbar}^{\mathrm{pp},w,M}(I)\bigr),
    d_M
  \right).
\]
The bracket, interaction, and differential are the transported
finite-window structures given by the admissibility datum.  The unweighted
Hamiltonian bracket enters only through that transport.\\
\hline
\(C_{n,M}\), \(C_{\Gamma,w}^{w_0}(\epsilon)\) &
Finite-window counterterms.  At order \(\hbar^n\),
\(C_{n,M}\in\mathcal Q^0_{w,M}(I)\) must satisfy
\[
  d_MC_{n,M}=-\mathfrak o^{\mathrm{ns}}_{n,M},
  \qquad
  \pi_{M,N}C_{n,M}=C_{n,N}.
\]
Graph representatives \(C_{\Gamma,w}^{w_0}(\epsilon)\) must also commute
with window truncation and weight transport.  They are degree-zero local
functionals in the filtered QME complex.\\
\hline
\(\mathfrak O^{\mathrm{ns}}_n\) &
Order-\(n\) finite-window obstruction vector
\[
  \mathfrak O^{\mathrm{ns}}_n
  =
  \bigl(([\mathfrak o^{\mathrm{ns}}_{n,M}])_M,\lambda_n\bigr),
\]
where the first component lies in
\(\varprojlim_MH^1(\mathcal Q^\bullet_{w,M}(I))\), and
\(\lambda_n\in\varprojlim\nolimits^1_M
H^0(\mathcal Q^\bullet_{w,M}(I))\) is the primitive-compatibility class.
Its vanishing is exactly the finite-window criterion for compatible
order-\(n\) counterterms.\\
\hline
\(\mathsf E_{\theta_3,M}\), \(C_{\theta_3,M}\) &
First proved larger-row non-scalar obstruction and its permitted
primitive.  In the one-row complex,
\[
  \mathcal Q^0_{\theta_3,M}=0,\qquad
  \mathcal Q^1_{\theta_3,M}=\C\mathsf E_{\theta_3,M},\qquad
  d_M=0,
\]
and the residual is \(\mathfrak o^{\mathrm{ns}}_{\theta_3,M}
=\mathsf E_{\theta_3,M}\).  The covector
\(\ell_{\theta_3}(\mathsf E_{\theta_3,M})=1\) proves that this class is
nonzero in the given subcomplex.  A valid \(C_{\theta_3,M}\) is one of
three explicit witnesses,
\[
  K^M_{\Theta_3}(\eta_{\theta_3,M}),\qquad
  C^{\mathrm{ct}}_{\theta_3,M},\qquad
  \text{or a companion-face primitive,}
\]
with differential \(d_MC_{\theta_3,M}=-\mathsf E_{\theta_3,M}\),
scalar-contact value zero, and vanishing Roos compatibility class under
finite-window projection.\\
\hline
\(I_N,\eta,\chi_N\) &
Moment-ideal and scalar primitive data.  The quantum Hamiltonian
reduction uses the left moment-map ideal
\(I_N=\mathcal W_N\widehat\mu(\lie{sl}_N)\) and then removes the scalar
empty trace.  The scalar primitive lives on the unreduced Hamiltonian
source \(A=\C[z_1,z_2]\):
\[
  \eta(f)=-[f]_0,\qquad d_{\mathrm{CE}}\eta=\omega,\qquad
  \chi_N(f)=N[f]_0=-N\eta(f).
\]
Thus
\(\hbar N\omega+d_{\mathrm{CE}}(\hbar\chi_N)=0\) before scalar
reduction.  This primitive lives on \(A\); after passage to
\(\bar A=A/\C\cdot1\), the scalar-reduced QME complex requires its own
counterterm data.\\
\hline
\(E_{a,b,N}\), \(A^j{}_i(a,b,N)\) &
Quantum shear moment-ideal primitive datum in the matrix Weyl algebra.
Put
\[
  E_{a,b,N}
  =
  \sum_{j=0}^b\operatorname{Tr}(Y^jX^aY^{b-j})
  -
  J_N\!\left(
    \{z_1^{a+1}/(a+1),z_2^{b+1}\}_\hbar
  \right).
\]
The all-\(N\) shear congruence required for the radial/Weyl image is
equivalent to
\[
  E_{a,b,N}\in I_N
  \quad\Longleftrightarrow\quad
  E_{a,b,N}
    =
  \sum_{i,j}A^j{}_i(a,b,N)\widehat\mu^i{}_j
\]
for coefficients \(A^j{}_i(a,b,N)\in\mathcal W_N\).  This is a
degree-zero left-ideal primitive problem in the quantum moment-map ideal.
The scalar CE primitive \(\eta\), scalar-contact projection, and Costello
QME counterterms are independent finite-row data.\\
\hline
Reduced principal-part brackets &
Target-side brackets in
\(A_{\partial,\hbar}^{\mathrm{pp},w}(I)\) after scalar reduction and
de Rham-current contraction:
\[
  \{B_f(a),B_g(b)\}_\hbar=B_{\{f,g\}_\hbar}(ab),\quad
  \{B_f(a),\Theta_\rho(B)\}_\hbar
    =\Theta_{\operatorname{ad}^{\vee}_{f,\hbar}\rho}(aB),
\]
\[
  \{\Theta_\rho(B),\Theta_\sigma(C)\}_\hbar=0 .
\]
They use multiplication of a distribution by a smooth test function.
They prove the reduced current interface.  Costello wavefront estimates,
unreduced centrality homotopies, and the non-scalar QME are additional
analytic and homotopical data.\\
\hline
Unreduced non-scalar QME realization &
The full graph theorem is the construction of
\(\kappa_{\hbar,w,I}\) and compatible counterterms
\(C_{\hbar,w}\).  Put
\(K_{\hbar,w,I}=\kappa_{\hbar,w,I}+C_{\hbar,w}\).  The equation is
\[
  \operatorname{Curv}(K_{\hbar,w,I})
  =
  d_{\mathrm{QME}}K_{\hbar,w,I}
  +\frac12[K_{\hbar,w,I},K_{\hbar,w,I}]_\hbar
  =0.
\]
It requires the Costello finite-scale BV propagator and Laplacian,
locality and counterterm estimates, scalar-contact chain projection,
vanishing of \(\mathfrak O^{\mathrm{ns}}_n\) for every order, and
unreduced product and \(P_0\)-bracket homotopies against arbitrary open
observables.\\
\hline
\end{longtable}

\paragraph{Common local rings and quotients.}
\(\mathfrak h\) labels Hamiltonian observables;
\(\mathfrak g_{\mathrm{str}}\) labels CE/PV closed cochains;
\(\bar A\) labels the polynomial trace sector; \(N\) labels the Dirac
brane stack rank.
\begin{center}
\begin{tabular}{@{}>{\raggedright\arraybackslash}p{0.25\textwidth}
  >{\raggedright\arraybackslash}p{0.67\textwidth}@{}}
\hline
\textbf{Symbol} & \textbf{Meaning}\\
\hline
\(R=\C[[z_1,z_2]]\) &
Complete local ring of the formal symplectic polydisk.\\
\hline
\(\mathfrak m=(z_1,z_2)\) &
Maximal ideal of functions vanishing at the brane point.\\
\hline
\(\mathfrak h=R/\C\cdot1\) &
Reduced Hamiltonian Lie algebra of the formal polydisk. Constants are
removed because they generate the zero Hamiltonian vector field.\\
\hline
\(\mathfrak g_{\mathrm{str}}\) &
Strict shifted cotangent Hamiltonian Lie algebra
\[
  \mathfrak h\ltimes\mathfrak h^\vee_{\mathrm{cont}}[1].
\]
Its Hamiltonian side is a product completion and its cotangent side is
continuous dual/direct-sum in residue coordinates.  This is the
formal-local CE/PV object.  A Costello coefficient kernel additionally
requires the endpoint Casimir data.\\
\hline
\(\mathfrak g_w\) &
Weighted finite-window Tate replacement used as a regulator.  It carries
the same coordinate-window CE/PV formulas after passage through the
admissible tower.  Its weight topology is auxiliary data for the
admissible tower; the strict cotangent completion is
\(\mathfrak g_{\mathrm{str}}\).\\
\hline
\(\bar A=\C[z_1,z_2]/\C\cdot1\) &
Polynomial Hamiltonian sector used in explicit trace and Moyal
calculations.\\
\hline
\(\bar A_{\mathrm{desc}}\) &
Abstract polynomial label space for PBW one-\(\psi\) descendants; it
carries a different \(\mathfrak h\)-action from the Matlis
principal-part module \(\mathcal P\) at the same bidegree labels.\\
\hline
\([\bar c]\) &
Distinguished scalar \(U(1)\) anomaly class in
\(H^2_{\mathrm{Lie}}(\bar A,\C)\), realized by the constant term of the
Poisson bracket and by the Capelli shift.\\
\hline
\(N\) &
Rank of the Dirac brane stack. Degreewise stable trace statements are
formed at fixed total word degree and then for \(N\) in the stable range.\\
\hline
\(D\) &
Taylor truncation degree when a finite quotient
\(\mathfrak m/\mathfrak m^{D+1}\) is used.  The brane rank is denoted
\(N\).\\
\hline
\end{tabular}
\end{center}

\paragraph{\texorpdfstring{Large-\(N\) and quotient conventions.}{Large-N and quotient conventions}}
An algebraic primitive projection is defined inside the stable range;
a physical cumulant requires a state \(\omega_{N,\Omega}\).
\begin{center}
\begin{tabular}{@{}>{\raggedright\arraybackslash}p{0.25\textwidth}
  >{\raggedright\arraybackslash}p{0.67\textwidth}@{}}
\hline
\textbf{Symbol or operation} & \textbf{Meaning}\\
\hline
\(\operatorname{Prim}_{\mathrm{st}}\) &
Degreewise stable primitive invariant sector. In total word degree
\(d\), the Procesi--Razmyslov stable range is \(N\geq d\).\\
\hline
\(\operatorname{Conn}_{N\to\infty}\) &
Connected primitive/indecomposable projection in the algebraic stable
trace sector after finite-\(N\) normalization and passage to the stable
range. Physical cumulants require a state \(\omega_{N,\Omega}\).\\
\hline
Hamiltonian scalar quotient &
The closed scalar \(1\in\C[z_1,z_2]\) and the open empty trace
\(\operatorname{Tr}_N(1)=N\) are omitted as primitive Hamiltonian
generators. The unit of the symmetric algebra remains.\\
\hline
\end{tabular}
\end{center}

\paragraph{Separation rule.}
PBW special-fibre labels \(H_{a,b},O_{a,b},\Psi_{a,b}\), CE/PV labels
\(c^I,u_I,\theta^I,O_I\), and Matlis labels \(\rho_{a,b}\) carry
different \(\mathfrak h\)-actions and different polarizations.  Shared
indices are labels for separate modules.

The Fourier--Rees theorem
(Theorem~\ref{thm:app-matlis-rees-fourier-bridge}) is the only
comparison morphism between PBW and Matlis labels: it matches
associated-graded labels through the Fourier delta Rees lattice
\(\mathcal D_\hbar/(z_1,z_2)\), with \v{C}ech realization carrying the
factor \((-1)^{a+b}\hbar^{a+b}a!b!\).  Before \(\hbar\) is inverted
this is a Rees sublattice; as a module statement it uses a
Fourier-twisted polarization.  The Hamiltonian coadjoint action is the
action on \(\mathcal P\).  Beyond the Fourier--Rees theorem the spans of
\(\Psi_{a,b}\) and \(\rho_{a,b}\) carry distinct
\(\mathfrak h\)-module structures.

\paragraph{Source--comparison table.}

\begin{longtable}{@{}>{\raggedright\arraybackslash}p{0.22\textwidth}
  >{\raggedright\arraybackslash}p{0.22\textwidth}
  >{\raggedright\arraybackslash}p{0.18\textwidth}
  >{\raggedright\arraybackslash}p{0.28\textwidth}@{}}
\hline
\textbf{Source (algebraic)} & \textbf{Target (physical)} &
\textbf{Comparison morphism} &
\textbf{Scope}\\
\hline
\endfirsthead
\hline
\textbf{Source (algebraic)} & \textbf{Target (physical)} &
\textbf{Comparison morphism} &
\textbf{Scope}\\
\hline
\endhead
coordinate-window tower / mixed-continuous replacement attached to
\(\mathfrak g_{\mathrm{str}}\) &
\(\mathrm{PV}_{\mathrm{coord}}(A_{\partial,\mathrm{coord}})\) &
coordinate dg identity; bracketed comparison only after choosing a
bracket-admissible \(B\) &
full unreduced factorization centre requires centrality and descent data\\
\hline
bracket-admissible \(B\subset\mathrm{PV}_{\mathrm{coord}}\) &
\(P_0\) realization \(\Phi^B_{P_0}\) &
restriction of \(\Phi_{\mathrm{coord}}\) when \(B\) gives completeness,
\(d_\pi^{\mathrm{coord}}\)-stability, product stability, and a
continuous contraction bracket &
bracket-admissible subalgebra inside the chosen completion\\
\hline
kernel-admissible \(B\) &
Maurer--Cartan kernel \(\Theta_B\) &
diagonal-convergent kernel in the chosen completed tensor product &
strict ambient product completion requires separate Casimir convergence\\
\hline
\(C^\bullet_{\mathrm{CE}}(\mathfrak g)\) / \(\mathrm{PV}(A)\) &
reduced \(P_0\) central-operation model &
local Poisson cochain model &
unreduced \(Z^{P_0}_{\mathrm{fact}}\) requires centrality homotopies\\
\hline
\(U_\hbar(\mathfrak g)/(\hbar)\) PBW special fibre &
\(\mathbf S(\mathfrak h\oplus\mathfrak h^\vee[1])\) and then selected
\(\mathbf S(\bar A\oplus\bar A_{\mathrm{desc}}[1])\) labels &
associated-graded label comparison &
PBW and Matlis labels carry different actions\\
\hline
admissible Tate / relative bar-cobar envelope &
\(\widehat U(\mathfrak g_{\mathrm{adm}})\) &
PBW envelope under conilpotent CE chains, exact continuous duality,
PBW completeness, and strong convergence &
available after specifying admissibility data for the formal polydisk\\
\hline
\(c^I\) &
\(\theta^I\) &
generator rule \(c^I\mapsto\theta^I\) of \(\Phi_{\mathrm{coord}}\) &
CE ghost coordinate\\
\hline
\(u_I\) &
\(O_I\) and brane-line Hamiltonian coordinates \(B_f(a)\) &
generator rule \(u_I\mapsto O_I\) of \(\Phi_{\mathrm{coord}}\) &
boundary Hamiltonian coordinate\\
\hline
\(C^\bullet_{\mathrm{CE}}(\mathfrak g_I)\) on intervals &
reduced brane-line factorization target &
cofactorization/\allowbreak{}opposite-interval source; factorization only after the
CE coproduct is paired with target multiplication &
variance reversal handled by the CE coproduct and target multiplication\\
\hline
\(\Psi_{a,b}\) &
stable primitive one-\(\psi\) PBW special-fibre class &
stable-trace identification in the trace complex &
PBW trace label with polynomial action\\
\hline
\(\rho_{a,b}\) &
\(\mathcal P=\operatorname{Ann}_{\mathrm{res}}(\C\cdot1)\) &
residue/Matlis pairing
\(f\otimes\rho\mapsto\operatorname{Res}_0(f\rho)\) &
Matlis principal-part label with Hamiltonian coadjoint action\\
\hline
\(\mathsf F_{\mathrm{Rees}}\) &
\(\bigoplus\hbar^{a+b}\C[[\hbar]]\rho_{a,b}\) &
Fourier-polarized Rees label bridge &
polarization-changing Rees comparison\\
\hline
\(\Omega_c^\bullet(I)\otimes\mathfrak h\) &
brane-line Hamiltonian currents &
reduced current model &
\(C_c^\infty(I)\) denotes only the degree-zero summand\\
\hline
\(\mathcal D'_c(I)\widehat\otimes\mathcal P[1]\) &
principal-part current sector \(\Theta_\rho(B)\) &
post-de-Rham-current contraction &
uses multiplication of a distribution by a smooth test function\\
\hline
formal Moyal algebra \(\bar A_\hbar\) &
degree-zero Weyl coefficient target &
reduced algebraic Moyal map with raw weights \(1\) and \(1/24\) &
finite-\(N\) radial-parts theorem requires the radial image data\\
\hline
finite-\(N\) radial-parts target &
stable connected trace realization of the Moyal coefficients &
finite verified relative datum: formal coefficients, Capelli
normalization, free normal-diagram obstruction
\(\mathfrak o^{\mathrm{rad}}_{a,b}\), edge PBW strip, and the
remaining uniform splitting or left-cokernel theorem &
requires formal Moyal coefficients, the radial-parts hypothesis, and the
normal-diagram obstruction calculation; Costello analytic graph/QME
realization is a separate construction\\
\hline
Costello brane-defect kernel &
closed-open BV/QME specialization &
obstruction datum
\(\operatorname{Ob}_{w,\partial,\hbar}\) of bulk-to-defect curvature &
requires bulk-to-defect curvature and counterterm data\\
\hline
scalar cocycle \([\bar c]\) &
\(U(1)\) anomaly / Capelli contact line &
distinguished class in
\(H^2_{\mathrm{Lie}}(\bar A,\C)\) &
one scalar contact class inside the unreduced QME problem\\
\hline
\end{longtable}

\paragraph{PBW special-fibre open trace labels.}
Algebraic side: the formal-stalk chart at the \(\hbar=0\) special fibre.
Physical side: the trace observable
\(\overline{\operatorname{Tr}\,f(\phi_1,\phi_2)}\) on the matrix Weyl
brane after moment-map reduction.  Comparison morphism: the PBW limit
of that trace observable, an isomorphism of stable trace complexes in
the Procesi--Razmyslov stable range \(N\geq d\).
\begin{center}
\begin{tabular}{@{}>{\raggedright\arraybackslash}p{0.25\textwidth}
  >{\raggedright\arraybackslash}p{0.67\textwidth}@{}}
\hline
\textbf{Symbol} & \textbf{Meaning}\\
\hline
\(H_{a,b}=z_1^az_2^b\) &
Polynomial Hamiltonian label in \(\bar A\), with \(a+b>0\). In the
PBW special fibre it labels the scalar-reduced trace observable
\(O_{a,b}\).\\
\hline
\(O_{a,b}\) &
Open ghost-zero trace class
\(\overline{\operatorname{Tr}(\phi_1^a\phi_2^b)}\), with \(a+b>0\)
in the reduced Hamiltonian label set.\\
\hline
\(\Psi_{a,b}\) &
Primitive one-\(\psi\) stable open class, represented by the
symmetrised trace with one \(\psi\), \(a\) copies of \(\phi_1\), and
\(b\) copies of \(\phi_2\), with \(a+b>0\) after removing the scalar
descendant line. It lies in the PBW special-fibre open
shadow and carries the corresponding polynomial action. In the reduced
comparison it gives the primitive indecomposable \(P_0\)-shadow
label.\\
\hline
\(\mathsf F_{\mathrm{Rees}}\) &
Fourier--Rees label bridge of
Theorem~\ref{thm:app-matlis-rees-fourier-bridge}. It sends the
polynomial label \(z_1^az_2^b\) to the Fourier-delta label
\(\partial_1^a\partial_2^be_0\), whose \v{C}ech realization is
\((-1)^{a+b}\hbar^{a+b}a!b!\rho_{a,b}\).  After inverting
\(\hbar\), division by this Rees factor gives the
principal-part label \(\rho_{a,b}\).  It changes polarization and is
the comparison between the PBW label \(\Psi_{a,b}\) and the Matlis label
\(\rho_{a,b}\).\\
\hline
\end{tabular}
\end{center}

The shared indices on \(H_{a,b}\), \(\Psi_{a,b}\), and
\(\rho_{a,b}\) are labels for separate modules: the span of the
\(\Psi_{a,b}\) is the PBW special-fibre trace label space, while
\(\mathfrak h^\vee_{\mathrm{cont}}\cong\mathcal P\) is the
principal-part current module.  The Fourier--Rees bridge
\(\mathsf F_{\mathrm{Rees}}\) is the only comparison morphism between
PBW and Matlis sectors.

\paragraph{Closed CE/PV labels.}
Algebraic side: the closed-side complex of CE/PV cochains.  Physical
side: the closed-string \(P_0\)-central operation on brane
observables.  Comparison morphism: \(\Phi_{\mathrm{coord}}\), the
coordinate dg CE/PV isomorphism on generators
\(c^I\mapsto\theta^I\), \(u_I\mapsto O_I\); its restriction to a
bracket-admissible \(B\) is the \(P_0\) realization \(\Phi^B_{P_0}\).
The parities are
\[
  |c^I|=|\theta^I|=1,\qquad |u_I|=|O_I|=0,
\]
and the cotangent / Schouten bracket has degree \(-1\), with
\([\theta^I,O_J]_S=\delta^I_J\).
\begin{center}
\begin{tabular}{@{}>{\raggedright\arraybackslash}p{0.25\textwidth}
  >{\raggedright\arraybackslash}p{0.67\textwidth}@{}}
\hline
\textbf{Symbol} & \textbf{Meaning}\\
\hline
\(\mathfrak g=\mathfrak h\ltimes
\mathfrak h^\vee_{\mathrm{cont}}[1]\) &
Shifted cotangent Lie algebra of Hamiltonian canonical transformations.\\
\hline
\(H_I\) &
Hamiltonian generator in \(\mathfrak h\).\\
\hline
\(\eta^I\) &
Shifted cotangent generator \((H_I)^\vee[1]\), with \(I=(a,b)\) and
\(a+b>0\).\\
\hline
\(c^I\) &
Chevalley--Eilenberg coordinate dual to a Hamiltonian generator
\(H_I\in\mathfrak h\). Under \(\Phi_{\mathrm{coord}}\),
\(c^I\mapsto\theta^I\).\\
\hline
\(u_I\) &
Chevalley--Eilenberg coordinate dual to the shifted cotangent generator.
Under \(\Phi_{\mathrm{coord}}\), \(u_I\mapsto O_I\). Boundary Hamiltonian observables
come from \(u_I\).\\
\hline
\(\theta^I\) &
Polyvector coordinate corresponding to \(c^I\) under the coordinate CE/PV
isomorphism.\\
\hline
\(O_I\) &
Hamiltonian observable coordinate corresponding to \(u_I\) under the
coordinate CE/PV isomorphism.\\
\hline
\(\Phi_{\mathrm{coord}}\) &
Coordinate dg CE/PV map determined on generators by
\(c^I\mapsto\theta^I\) and \(u_I\mapsto O_I\).  Its \(P_0\) restriction
is \(\Phi^B_{P_0}\) only after a bracket-admissible \(B\) is chosen.\\
\hline
\end{tabular}
\end{center}

\paragraph{Matlis principal-part labels.}
Algebraic side: the Grothendieck-dual of the formal jet algebra
\(\mathfrak h\), namely the residue/Matlis principal-part module
\(\mathcal P\cong\mathfrak h^\vee_{\mathrm{cont}}\).  Physical side:
the principal-part current observable \(\Theta_\rho(B)\) on a brane
interval.  Comparison morphism: the residue pairing
\(\mathfrak h\otimes\mathcal P\to\C\),
\(f\otimes\rho\mapsto\operatorname{Res}_0(f\rho)\), with Hamiltonian
coadjoint action on \(\mathcal P\) its image.  This is the
formal-holomorphic locality (iii) of
Theorem~\ref{thm:dict-kinematic-locality}.
\begin{center}
\begin{tabular}{@{}>{\raggedright\arraybackslash}p{0.25\textwidth}
  >{\raggedright\arraybackslash}p{0.67\textwidth}@{}}
\hline
\textbf{Symbol} & \textbf{Meaning}\\
\hline
\(\mathcal P\) &
Reduced principal-part module
\(\operatorname{Ann}_{\mathrm{res}}(\C\cdot1)\subset
H^2_{\mathfrak m}(R)\,dz_1dz_2\), identified with
\(\mathfrak h^\vee_{\mathrm{cont}}\).\\
\hline
\(\rho_{a,b}\) &
Principal-part residue current
\(z_1^{-a-1}z_2^{-b-1}dz_1dz_2\), \(a+b>0\), with the Matlis
coadjoint action.\\
\hline
\(z_1^p z_2^q\cdot\rho_{a,b}\) &
Coadjoint action
\(-(pb-qa+p-q)\rho_{a-p+1,b-q+1}\). This is the Hamiltonian coadjoint
principal-part action.  The Matlis \(R\)-module action and the PBW action
on \(\Psi_{a,b}\) are separate structures.\\
\hline
\(\mathcal P_I\) &
\(\mathcal P_I=\mathcal D'_c(I)\widehat\otimes\mathcal P\), the
compactly supported distributional principal-part currents on an interval
\(I\).\\
\hline
\end{tabular}
\end{center}

As stated in the separation rule above,
\[
  \Psi_{a,b}\in\text{PBW special-fibre open homology},\qquad
  \rho_{a,b}\in
  \mathcal P\cong\mathfrak h^\vee_{\mathrm{cont}}
\]
live in different modules.  The only comparison morphism between them
is the Fourier--Rees label bridge \(\mathsf F_{\mathrm{Rees}}\) of
Theorem~\ref{thm:app-matlis-rees-fourier-bridge}, which acts on
associated-graded labels through Fourier-twisted polarization.  The
Hamiltonian coadjoint action remains the action on the Matlis module.

\paragraph{Interval-current and boundary labels.}
Algebraic side: the brane-line objects produced by the support-locality
differential of Theorem~\ref{thm:dict-kinematic-locality}\,(i).
Physical side: the smeared brane measurement \(B_f(a)\) and the
principal-part current observable \(\Theta_\rho(B)\).  Comparison
morphism: the de Rham-current contraction
\(u_{a(t)dt\otimes f}\mapsto B_f(a)\),
\(c_{B,\rho}\mapsto\Theta_\rho(B)\); a smooth test function multiplies
a distribution.
\begin{center}
\begin{tabular}{@{}>{\raggedright\arraybackslash}p{0.25\textwidth}
  >{\raggedright\arraybackslash}p{0.67\textwidth}@{}}
\hline
\textbf{Symbol} & \textbf{Meaning}\\
\hline
\(\mathfrak h_I\) &
\(\mathfrak h_I=\Omega^\bullet_c(I)\widehat\otimes\mathfrak h\), the
compactly supported de Rham Hamiltonian currents on an interval
\(I\subset\R\). Degree-zero test functions appear only after taking the
explicit measurement notation.\\
\hline
\(\mathfrak g_I\) &
\(\mathfrak g_I=\mathfrak h_I\ltimes\mathcal P_I[1]\), the shifted
cotangent current algebra on \(I\), with principal-part currents in the
shifted cotangent direction.\\
\hline
\(B_f(a)\) &
Smeared boundary Hamiltonian observable
\(\int_I a(t)\overline{\operatorname{Tr}f(\phi_1(t),\phi_2(t))}\,dt\).
Its CE source is \(u_{a(t)dt\otimes f}\).\\
\hline
\(A^{\mathrm{Ham}}_\partial(I)\) &
Reduced Hamiltonian brane-line \(P_0\) factorization algebra
\(\widehat{\mathbf S}(\mathfrak h_I)\).\\
\hline
\(A^{\mathrm{pp}}_\partial(I)\) &
Reduced principal-part defect-current model after de Rham-current contraction
\(\widehat{\mathbf S}(\mathfrak h_I)\widehat\otimes
\widehat{\mathbf S}(\mathcal P_I[1])\).  An unreduced BV observable
algebra maps to it in the presence of the centrality homotopies,
coefficient kernel, brane-defect propagator, descent data, and BV/QME
solution.\\
\hline
\(\Theta_\rho(B)\) &
Principal-part current observable with
\(B\in\mathcal D'_c(I)\) and \(\rho\in\mathcal P\). It is the boundary
notation for \(B\otimes\rho[1]\) in the reduced model and pairs with
smooth test functions by \(B(a)\operatorname{Res}(\rho f)\).\\
\hline
\(aB\) &
Multiplication of a compactly supported distribution \(B\) by a smooth
test function \(a\in\Omega^0_c(I)\). The reduced brackets use this
operation and never require multiplying two distributions.\\
\hline
\end{tabular}
\end{center}
\endgroup
